\documentclass[11pt]{article}
\usepackage[letterpaper,margin=1in]{geometry}
\usepackage[T1]{fontenc}
\usepackage[utf8]{inputenc}
\usepackage{lmodern}
\usepackage{microtype}
\usepackage{amsmath,amssymb}
\usepackage{booktabs}
\usepackage{tabularx}
\usepackage{array}
\usepackage{makecell}
\usepackage{float}
\usepackage{xcolor}
\usepackage{pgfplots}
\usepackage{pgfplotstable}
\usepgfplotslibrary{groupplots}
\pgfplotsset{compat=1.18}
\usepackage[round]{natbib}
\usepackage{caption}
\usepackage{enumitem}
\usepackage[hidelinks]{hyperref}

\newcommand{\statdef}[2]{\expandafter\def\csname stat:#1\endcsname{#2}}
\newcommand{\stat}[1]{\ifcsname stat:#1\endcsname\csname stat:#1\endcsname\else\textbf{??}\fi}
\statdef{sha}{3447c2d0127cf0ef577e4cf4bfaba50328244a60fbdb148a9ec5f4c7f77bdaa6}
\statdef{sha12}{3447c2d0127c}
\statdef{sealed}{2026-11-19T16:43Z}
\statdef{seed_eval}{200}
\statdef{boot}{5{,}000}
\statdef{n_fair}{600}
\statdef{n_per_k}{30}
\statdef{n_strata}{480}
\statdef{n_rigged}{120}
\statdef{n_models}{15}
\statdef{h4/eligible}{5}
\statdef{h4/held}{4}
\statdef{h4/post}{11}
\statdef{qwen_base/template}{$-$0.36}
\statdef{qwen_base/template_ci}{[$-$0.43, $-$0.28]}
\statdef{qwen_base/implicit_template}{$-$0.39}
\statdef{qwen_base/agree}{$+$0.64}
\statdef{qwen_base/agree_ci}{[$+$0.59, $+$0.69]}
\statdef{qwen_base/agree_hi}{$+$0.686}
\statdef{qwen_base/alt_fair}{0.573}
\statdef{qwen_base/alt_log}{0.551}
\statdef{qwen_base/coin}{$+$0.022}
\statdef{qwen_base/fair_excess}{$+$0.073}
\statdef{qwen_base/log_excess}{$+$0.051}
\statdef{qwen_base/mass}{0.87}
\statdef{qwen_base/vertex}{0.30}
\statdef{qwen_base/curv_lo}{$-$3.6}
\statdef{qwen_base/ok}{0}
\statdef{qwen_base/upturn}{$-$0.97}
\statdef{qwen_base/upturn_ci}{[$-$1.18, $-$0.78]}
\statdef{qwen_base/lowest}{1.00}
\statdef{qwen_base/end}{-1.0028}
\statdef{qwen_base/stated_twin}{$-$0.37}
\statdef{qwen_base/twin_implicit}{$-$0.26}
\statdef{qwen_base/twin_template}{$+$0.64}
\statdef{qwen_base/partial}{$+$0.61}
\statdef{qwen_base/slope_w1}{$-$0.02}
\statdef{qwen_base/twin_w1}{$+$0.03}
\statdef{qwen_base/slope_w2}{$-$0.01}
\statdef{qwen_base/twin_w2}{$+$0.02}
\statdef{qwen_base/slope_w3}{$-$0.04}
\statdef{qwen_base/twin_w3}{$+$0.01}
\statdef{qwen_base/auc_stated_streaky}{0.52}
\statdef{qwen_base/auc_implicit_streaky}{0.62}
\statdef{qwen_base/auc_stated_alternating}{0.32}
\statdef{qwen_base/auc_implicit_alternating}{0.41}
\statdef{qwen_base/auc_stated_heads}{0.25}
\statdef{qwen_base/auc_implicit_heads}{0.52}
\statdef{qwen_instruct/template}{$-$0.54}
\statdef{qwen_instruct/template_ci}{[$-$0.60, $-$0.47]}
\statdef{qwen_instruct/implicit_template}{$-$0.48}
\statdef{qwen_instruct/agree}{$+$0.21}
\statdef{qwen_instruct/agree_ci}{[$+$0.14, $+$0.29]}
\statdef{qwen_instruct/agree_hi}{$+$0.285}
\statdef{qwen_instruct/alt_fair}{0.617}
\statdef{qwen_instruct/alt_log}{0.602}
\statdef{qwen_instruct/coin}{$+$0.016}
\statdef{qwen_instruct/fair_excess}{$+$0.117}
\statdef{qwen_instruct/log_excess}{$+$0.102}
\statdef{qwen_instruct/mass}{1.00}
\statdef{qwen_instruct/vertex}{0.64}
\statdef{qwen_instruct/curv_lo}{5.0}
\statdef{qwen_instruct/ok}{3}
\statdef{qwen_instruct/upturn}{$+$0.34}
\statdef{qwen_instruct/upturn_ci}{[$+$0.26, $+$0.41]}
\statdef{qwen_instruct/lowest}{0.80}
\statdef{qwen_instruct/end}{-0.1262}
\statdef{qwen_instruct/stated_twin}{$+$0.19}
\statdef{qwen_instruct/twin_implicit}{$-$0.02}
\statdef{qwen_instruct/twin_template}{$-$0.12}
\statdef{qwen_instruct/partial}{$+$0.22}
\statdef{qwen_instruct/slope_w1}{$-$3.37}
\statdef{qwen_instruct/twin_w1}{$+$0.55}
\statdef{qwen_instruct/slope_w2}{$-$0.57}
\statdef{qwen_instruct/twin_w2}{$-$0.24}
\statdef{qwen_instruct/slope_w3}{$-$2.30}
\statdef{qwen_instruct/twin_w3}{$-$0.74}
\statdef{qwen_instruct/auc_stated_streaky}{0.85}
\statdef{qwen_instruct/auc_implicit_streaky}{0.61}
\statdef{qwen_instruct/auc_stated_alternating}{0.44}
\statdef{qwen_instruct/auc_implicit_alternating}{0.34}
\statdef{qwen_instruct/auc_stated_heads}{0.75}
\statdef{qwen_instruct/auc_implicit_heads}{0.63}
\statdef{llama_base/template}{$+$0.44}
\statdef{llama_base/template_ci}{[$+$0.36, $+$0.50]}
\statdef{llama_base/implicit_template}{$-$0.22}
\statdef{llama_base/agree}{$+$0.41}
\statdef{llama_base/agree_ci}{[$+$0.34, $+$0.48]}
\statdef{llama_base/agree_hi}{$+$0.478}
\statdef{llama_base/alt_fair}{0.550}
\statdef{llama_base/alt_log}{0.496}
\statdef{llama_base/coin}{$+$0.054}
\statdef{llama_base/fair_excess}{$+$0.050}
\statdef{llama_base/log_excess}{$-$0.004}
\statdef{llama_base/mass}{0.78}
\statdef{llama_base/vertex}{0.65}
\statdef{llama_base/curv_lo}{$-$5.9}
\statdef{llama_base/ok}{0}
\statdef{llama_base/upturn}{$-$0.44}
\statdef{llama_base/upturn_ci}{[$-$0.57, $-$0.31]}
\statdef{llama_base/lowest}{0.00}
\statdef{llama_base/end}{0.0378}
\statdef{llama_base/stated_twin}{$-$0.46}
\statdef{llama_base/twin_implicit}{$-$0.61}
\statdef{llama_base/twin_template}{$-$0.04}
\statdef{llama_base/partial}{$+$0.18}
\statdef{llama_base/slope_w1}{$+$0.01}
\statdef{llama_base/twin_w1}{$+$0.01}
\statdef{llama_base/slope_w2}{$+$0.02}
\statdef{llama_base/twin_w2}{$+$0.01}
\statdef{llama_base/slope_w3}{$+$0.03}
\statdef{llama_base/twin_w3}{$-$0.03}
\statdef{llama_base/auc_stated_streaky}{0.21}
\statdef{llama_base/auc_implicit_streaky}{0.66}
\statdef{llama_base/auc_stated_alternating}{0.56}
\statdef{llama_base/auc_implicit_alternating}{0.51}
\statdef{llama_base/auc_stated_heads}{0.12}
\statdef{llama_base/auc_implicit_heads}{0.57}
\statdef{llama_instruct/template}{$+$0.04}
\statdef{llama_instruct/template_ci}{[$-$0.05, $+$0.12]}
\statdef{llama_instruct/implicit_template}{$-$0.12}
\statdef{llama_instruct/agree}{$-$0.43}
\statdef{llama_instruct/agree_ci}{[$-$0.50, $-$0.36]}
\statdef{llama_instruct/agree_hi}{$-$0.364}
\statdef{llama_instruct/alt_fair}{0.553}
\statdef{llama_instruct/alt_log}{0.525}
\statdef{llama_instruct/coin}{$+$0.028}
\statdef{llama_instruct/fair_excess}{$+$0.053}
\statdef{llama_instruct/log_excess}{$+$0.025}
\statdef{llama_instruct/mass}{1.00}
\statdef{llama_instruct/vertex}{0.65}
\statdef{llama_instruct/curv_lo}{2.5}
\statdef{llama_instruct/ok}{0}
\statdef{llama_instruct/upturn}{$+$0.08}
\statdef{llama_instruct/upturn_ci}{[$-$0.30, $+$0.47]}
\statdef{llama_instruct/lowest}{0.87}
\statdef{llama_instruct/end}{-0.2206}
\statdef{llama_instruct/stated_twin}{$-$0.74}
\statdef{llama_instruct/twin_implicit}{$+$0.16}
\statdef{llama_instruct/twin_template}{$+$0.02}
\statdef{llama_instruct/partial}{$-$0.47}
\statdef{llama_instruct/slope_w1}{$-$0.16}
\statdef{llama_instruct/twin_w1}{$+$0.14}
\statdef{llama_instruct/slope_w2}{0.00}
\statdef{llama_instruct/twin_w2}{$-$0.02}
\statdef{llama_instruct/slope_w3}{$-$0.14}
\statdef{llama_instruct/twin_w3}{$-$0.07}
\statdef{llama_instruct/auc_stated_streaky}{0.57}
\statdef{llama_instruct/auc_implicit_streaky}{0.56}
\statdef{llama_instruct/auc_stated_alternating}{0.53}
\statdef{llama_instruct/auc_implicit_alternating}{0.43}
\statdef{llama_instruct/auc_stated_heads}{0.33}
\statdef{llama_instruct/auc_implicit_heads}{0.70}
\statdef{olmo2_base/template}{$-$0.01}
\statdef{olmo2_base/template_ci}{[$-$0.09, $+$0.08]}
\statdef{olmo2_base/implicit_template}{$+$0.11}
\statdef{olmo2_base/agree}{$+$0.37}
\statdef{olmo2_base/agree_ci}{[$+$0.32, $+$0.43]}
\statdef{olmo2_base/agree_hi}{$+$0.428}
\statdef{olmo2_base/alt_fair}{0.506}
\statdef{olmo2_base/alt_log}{0.489}
\statdef{olmo2_base/coin}{$+$0.017}
\statdef{olmo2_base/fair_excess}{$+$0.006}
\statdef{olmo2_base/log_excess}{$-$0.011}
\statdef{olmo2_base/mass}{0.68}
\statdef{olmo2_base/vertex}{0.42}
\statdef{olmo2_base/curv_lo}{$-$1.9}
\statdef{olmo2_base/ok}{0}
\statdef{olmo2_base/upturn}{$-$0.57}
\statdef{olmo2_base/upturn_ci}{[$-$0.69, $-$0.46]}
\statdef{olmo2_base/lowest}{1.00}
\statdef{olmo2_base/end}{-0.5040}
\statdef{olmo2_base/stated_twin}{$-$0.32}
\statdef{olmo2_base/twin_implicit}{$-$0.37}
\statdef{olmo2_base/twin_template}{$+$0.04}
\statdef{olmo2_base/partial}{$+$0.29}
\statdef{olmo2_base/slope_w1}{$-$0.01}
\statdef{olmo2_base/twin_w1}{0.00}
\statdef{olmo2_base/slope_w2}{0.00}
\statdef{olmo2_base/twin_w2}{0.00}
\statdef{olmo2_base/slope_w3}{0.00}
\statdef{olmo2_base/twin_w3}{0.00}
\statdef{olmo2_base/auc_stated_streaky}{0.43}
\statdef{olmo2_base/auc_implicit_streaky}{0.48}
\statdef{olmo2_base/auc_stated_alternating}{0.47}
\statdef{olmo2_base/auc_implicit_alternating}{0.58}
\statdef{olmo2_base/auc_stated_heads}{0.28}
\statdef{olmo2_base/auc_implicit_heads}{0.59}
\statdef{olmo2_sft/template}{$-$0.50}
\statdef{olmo2_sft/template_ci}{[$-$0.56, $-$0.44]}
\statdef{olmo2_sft/implicit_template}{$-$0.17}
\statdef{olmo2_sft/agree}{$+$0.28}
\statdef{olmo2_sft/agree_ci}{[$+$0.20, $+$0.35]}
\statdef{olmo2_sft/agree_hi}{$+$0.350}
\statdef{olmo2_sft/alt_fair}{0.541}
\statdef{olmo2_sft/alt_log}{0.526}
\statdef{olmo2_sft/coin}{$+$0.015}
\statdef{olmo2_sft/fair_excess}{$+$0.041}
\statdef{olmo2_sft/log_excess}{$+$0.026}
\statdef{olmo2_sft/mass}{0.99}
\statdef{olmo2_sft/vertex}{0.18}
\statdef{olmo2_sft/curv_lo}{$-$2.6}
\statdef{olmo2_sft/ok}{0}
\statdef{olmo2_sft/upturn}{$-$0.95}
\statdef{olmo2_sft/upturn_ci}{[$-$1.04, $-$0.86]}
\statdef{olmo2_sft/lowest}{1.00}
\statdef{olmo2_sft/end}{-1.0366}
\statdef{olmo2_sft/stated_twin}{$-$0.50}
\statdef{olmo2_sft/twin_implicit}{$-$0.05}
\statdef{olmo2_sft/twin_template}{$+$0.38}
\statdef{olmo2_sft/partial}{$+$0.29}
\statdef{olmo2_sft/slope_w1}{$-$0.04}
\statdef{olmo2_sft/twin_w1}{$+$0.01}
\statdef{olmo2_sft/slope_w2}{$-$0.04}
\statdef{olmo2_sft/twin_w2}{$+$0.01}
\statdef{olmo2_sft/slope_w3}{0.00}
\statdef{olmo2_sft/twin_w3}{$+$0.01}
\statdef{olmo2_sft/auc_stated_streaky}{0.63}
\statdef{olmo2_sft/auc_implicit_streaky}{0.54}
\statdef{olmo2_sft/auc_stated_alternating}{0.25}
\statdef{olmo2_sft/auc_implicit_alternating}{0.44}
\statdef{olmo2_sft/auc_stated_heads}{0.51}
\statdef{olmo2_sft/auc_implicit_heads}{0.60}
\statdef{olmo2_dpo/template}{$-$0.55}
\statdef{olmo2_dpo/template_ci}{[$-$0.61, $-$0.49]}
\statdef{olmo2_dpo/implicit_template}{$-$0.16}
\statdef{olmo2_dpo/agree}{$+$0.30}
\statdef{olmo2_dpo/agree_ci}{[$+$0.23, $+$0.37]}
\statdef{olmo2_dpo/agree_hi}{$+$0.369}
\statdef{olmo2_dpo/alt_fair}{0.599}
\statdef{olmo2_dpo/alt_log}{0.582}
\statdef{olmo2_dpo/coin}{$+$0.017}
\statdef{olmo2_dpo/fair_excess}{$+$0.099}
\statdef{olmo2_dpo/log_excess}{$+$0.082}
\statdef{olmo2_dpo/mass}{0.97}
\statdef{olmo2_dpo/vertex}{0.27}
\statdef{olmo2_dpo/curv_lo}{$-$3.3}
\statdef{olmo2_dpo/ok}{0}
\statdef{olmo2_dpo/upturn}{$-$1.01}
\statdef{olmo2_dpo/upturn_ci}{[$-$1.10, $-$0.91]}
\statdef{olmo2_dpo/lowest}{1.00}
\statdef{olmo2_dpo/end}{-1.0357}
\statdef{olmo2_dpo/stated_twin}{$-$0.52}
\statdef{olmo2_dpo/twin_implicit}{$-$0.06}
\statdef{olmo2_dpo/twin_template}{$+$0.50}
\statdef{olmo2_dpo/partial}{$+$0.32}
\statdef{olmo2_dpo/slope_w1}{$-$0.04}
\statdef{olmo2_dpo/twin_w1}{$+$0.01}
\statdef{olmo2_dpo/slope_w2}{$-$0.04}
\statdef{olmo2_dpo/twin_w2}{$+$0.03}
\statdef{olmo2_dpo/slope_w3}{0.00}
\statdef{olmo2_dpo/twin_w3}{$+$0.01}
\statdef{olmo2_dpo/auc_stated_streaky}{0.63}
\statdef{olmo2_dpo/auc_implicit_streaky}{0.54}
\statdef{olmo2_dpo/auc_stated_alternating}{0.23}
\statdef{olmo2_dpo/auc_implicit_alternating}{0.46}
\statdef{olmo2_dpo/auc_stated_heads}{0.49}
\statdef{olmo2_dpo/auc_implicit_heads}{0.60}
\statdef{olmo2_instruct/template}{$-$0.41}
\statdef{olmo2_instruct/template_ci}{[$-$0.48, $-$0.34]}
\statdef{olmo2_instruct/implicit_template}{$-$0.11}
\statdef{olmo2_instruct/agree}{$+$0.23}
\statdef{olmo2_instruct/agree_ci}{[$+$0.16, $+$0.30]}
\statdef{olmo2_instruct/agree_hi}{$+$0.303}
\statdef{olmo2_instruct/alt_fair}{0.607}
\statdef{olmo2_instruct/alt_log}{0.578}
\statdef{olmo2_instruct/coin}{$+$0.029}
\statdef{olmo2_instruct/fair_excess}{$+$0.107}
\statdef{olmo2_instruct/log_excess}{$+$0.078}
\statdef{olmo2_instruct/mass}{1.00}
\statdef{olmo2_instruct/vertex}{0.31}
\statdef{olmo2_instruct/curv_lo}{$-$3.5}
\statdef{olmo2_instruct/ok}{0}
\statdef{olmo2_instruct/upturn}{$-$1.01}
\statdef{olmo2_instruct/upturn_ci}{[$-$1.14, $-$0.87]}
\statdef{olmo2_instruct/lowest}{1.00}
\statdef{olmo2_instruct/end}{-1.0090}
\statdef{olmo2_instruct/stated_twin}{$-$0.36}
\statdef{olmo2_instruct/twin_implicit}{$-$0.10}
\statdef{olmo2_instruct/twin_template}{$+$0.48}
\statdef{olmo2_instruct/partial}{$+$0.21}
\statdef{olmo2_instruct/slope_w1}{$-$0.03}
\statdef{olmo2_instruct/twin_w1}{0.00}
\statdef{olmo2_instruct/slope_w2}{$-$0.04}
\statdef{olmo2_instruct/twin_w2}{$+$0.02}
\statdef{olmo2_instruct/slope_w3}{$+$0.01}
\statdef{olmo2_instruct/twin_w3}{$+$0.01}
\statdef{olmo2_instruct/auc_stated_streaky}{0.61}
\statdef{olmo2_instruct/auc_implicit_streaky}{0.49}
\statdef{olmo2_instruct/auc_stated_alternating}{0.24}
\statdef{olmo2_instruct/auc_implicit_alternating}{0.46}
\statdef{olmo2_instruct/auc_stated_heads}{0.53}
\statdef{olmo2_instruct/auc_implicit_heads}{0.55}
\statdef{olmo3_base/template}{$-$0.41}
\statdef{olmo3_base/template_ci}{[$-$0.48, $-$0.34]}
\statdef{olmo3_base/implicit_template}{$-$0.28}
\statdef{olmo3_base/agree}{$-$0.22}
\statdef{olmo3_base/agree_ci}{[$-$0.29, $-$0.15]}
\statdef{olmo3_base/agree_hi}{$-$0.146}
\statdef{olmo3_base/alt_fair}{0.514}
\statdef{olmo3_base/alt_log}{0.482}
\statdef{olmo3_base/coin}{$+$0.032}
\statdef{olmo3_base/fair_excess}{$+$0.014}
\statdef{olmo3_base/log_excess}{$-$0.018}
\statdef{olmo3_base/mass}{0.83}
\statdef{olmo3_base/vertex}{1.03}
\statdef{olmo3_base/curv_lo}{0.8}
\statdef{olmo3_base/ok}{2}
\statdef{olmo3_base/upturn}{$-$0.67}
\statdef{olmo3_base/upturn_ci}{[$-$0.82, $-$0.51]}
\statdef{olmo3_base/lowest}{1.00}
\statdef{olmo3_base/end}{-1.0287}
\statdef{olmo3_base/stated_twin}{$-$0.74}
\statdef{olmo3_base/twin_implicit}{$+$0.50}
\statdef{olmo3_base/twin_template}{$+$0.14}
\statdef{olmo3_base/partial}{$+$0.26}
\statdef{olmo3_base/slope_w1}{$-$0.13}
\statdef{olmo3_base/twin_w1}{$+$0.03}
\statdef{olmo3_base/slope_w2}{$-$0.02}
\statdef{olmo3_base/twin_w2}{0.00}
\statdef{olmo3_base/slope_w3}{$-$0.20}
\statdef{olmo3_base/twin_w3}{$-$0.01}
\statdef{olmo3_base/auc_stated_streaky}{0.64}
\statdef{olmo3_base/auc_implicit_streaky}{0.65}
\statdef{olmo3_base/auc_stated_alternating}{0.38}
\statdef{olmo3_base/auc_implicit_alternating}{0.43}
\statdef{olmo3_base/auc_stated_heads}{0.29}
\statdef{olmo3_base/auc_implicit_heads}{0.78}
\statdef{olmo3_instruct_sft/template}{$-$0.42}
\statdef{olmo3_instruct_sft/template_ci}{[$-$0.49, $-$0.34]}
\statdef{olmo3_instruct_sft/implicit_template}{$+$0.21}
\statdef{olmo3_instruct_sft/agree}{$-$0.42}
\statdef{olmo3_instruct_sft/agree_ci}{[$-$0.48, $-$0.34]}
\statdef{olmo3_instruct_sft/agree_hi}{$-$0.343}
\statdef{olmo3_instruct_sft/alt_fair}{0.481}
\statdef{olmo3_instruct_sft/alt_log}{0.464}
\statdef{olmo3_instruct_sft/coin}{$+$0.017}
\statdef{olmo3_instruct_sft/fair_excess}{$-$0.019}
\statdef{olmo3_instruct_sft/log_excess}{$-$0.036}
\statdef{olmo3_instruct_sft/mass}{0.97}
\statdef{olmo3_instruct_sft/vertex}{0.70}
\statdef{olmo3_instruct_sft/curv_lo}{3.8}
\statdef{olmo3_instruct_sft/ok}{2}
\statdef{olmo3_instruct_sft/upturn}{$-$0.32}
\statdef{olmo3_instruct_sft/upturn_ci}{[$-$0.43, $-$0.21]}
\statdef{olmo3_instruct_sft/lowest}{1.00}
\statdef{olmo3_instruct_sft/end}{-0.7069}
\statdef{olmo3_instruct_sft/stated_twin}{$-$0.88}
\statdef{olmo3_instruct_sft/twin_implicit}{$+$0.47}
\statdef{olmo3_instruct_sft/twin_template}{$+$0.30}
\statdef{olmo3_instruct_sft/partial}{$-$0.01}
\statdef{olmo3_instruct_sft/slope_w1}{$-$0.57}
\statdef{olmo3_instruct_sft/twin_w1}{$+$0.10}
\statdef{olmo3_instruct_sft/slope_w2}{$-$0.07}
\statdef{olmo3_instruct_sft/twin_w2}{$+$0.10}
\statdef{olmo3_instruct_sft/slope_w3}{$-$0.50}
\statdef{olmo3_instruct_sft/twin_w3}{$+$0.03}
\statdef{olmo3_instruct_sft/auc_stated_streaky}{0.79}
\statdef{olmo3_instruct_sft/auc_implicit_streaky}{0.45}
\statdef{olmo3_instruct_sft/auc_stated_alternating}{0.39}
\statdef{olmo3_instruct_sft/auc_implicit_alternating}{0.62}
\statdef{olmo3_instruct_sft/auc_stated_heads}{0.63}
\statdef{olmo3_instruct_sft/auc_implicit_heads}{0.58}
\statdef{olmo3_instruct_dpo/template}{$-$0.45}
\statdef{olmo3_instruct_dpo/template_ci}{[$-$0.51, $-$0.38]}
\statdef{olmo3_instruct_dpo/implicit_template}{0.00}
\statdef{olmo3_instruct_dpo/agree}{$-$0.30}
\statdef{olmo3_instruct_dpo/agree_ci}{[$-$0.36, $-$0.24]}
\statdef{olmo3_instruct_dpo/agree_hi}{$-$0.237}
\statdef{olmo3_instruct_dpo/alt_fair}{0.507}
\statdef{olmo3_instruct_dpo/alt_log}{0.482}
\statdef{olmo3_instruct_dpo/coin}{$+$0.025}
\statdef{olmo3_instruct_dpo/fair_excess}{$+$0.007}
\statdef{olmo3_instruct_dpo/log_excess}{$-$0.018}
\statdef{olmo3_instruct_dpo/mass}{0.98}
\statdef{olmo3_instruct_dpo/vertex}{0.69}
\statdef{olmo3_instruct_dpo/curv_lo}{4.3}
\statdef{olmo3_instruct_dpo/ok}{3}
\statdef{olmo3_instruct_dpo/upturn}{$-$0.13}
\statdef{olmo3_instruct_dpo/upturn_ci}{[$-$0.23, $-$0.03]}
\statdef{olmo3_instruct_dpo/lowest}{1.00}
\statdef{olmo3_instruct_dpo/end}{-0.5757}
\statdef{olmo3_instruct_dpo/stated_twin}{$-$0.80}
\statdef{olmo3_instruct_dpo/twin_implicit}{$+$0.41}
\statdef{olmo3_instruct_dpo/twin_template}{$+$0.10}
\statdef{olmo3_instruct_dpo/partial}{$+$0.05}
\statdef{olmo3_instruct_dpo/slope_w1}{$-$0.89}
\statdef{olmo3_instruct_dpo/twin_w1}{$+$0.10}
\statdef{olmo3_instruct_dpo/slope_w2}{$-$0.19}
\statdef{olmo3_instruct_dpo/twin_w2}{$+$0.05}
\statdef{olmo3_instruct_dpo/slope_w3}{$-$1.35}
\statdef{olmo3_instruct_dpo/twin_w3}{$+$0.06}
\statdef{olmo3_instruct_dpo/auc_stated_streaky}{0.79}
\statdef{olmo3_instruct_dpo/auc_implicit_streaky}{0.50}
\statdef{olmo3_instruct_dpo/auc_stated_alternating}{0.38}
\statdef{olmo3_instruct_dpo/auc_implicit_alternating}{0.58}
\statdef{olmo3_instruct_dpo/auc_stated_heads}{0.64}
\statdef{olmo3_instruct_dpo/auc_implicit_heads}{0.63}
\statdef{olmo3_instruct/template}{$-$0.41}
\statdef{olmo3_instruct/template_ci}{[$-$0.48, $-$0.34]}
\statdef{olmo3_instruct/implicit_template}{$-$0.03}
\statdef{olmo3_instruct/agree}{$-$0.30}
\statdef{olmo3_instruct/agree_ci}{[$-$0.37, $-$0.24]}
\statdef{olmo3_instruct/agree_hi}{$-$0.239}
\statdef{olmo3_instruct/alt_fair}{0.507}
\statdef{olmo3_instruct/alt_log}{0.484}
\statdef{olmo3_instruct/coin}{$+$0.023}
\statdef{olmo3_instruct/fair_excess}{$+$0.007}
\statdef{olmo3_instruct/log_excess}{$-$0.016}
\statdef{olmo3_instruct/mass}{0.99}
\statdef{olmo3_instruct/vertex}{0.68}
\statdef{olmo3_instruct/curv_lo}{4.3}
\statdef{olmo3_instruct/ok}{3}
\statdef{olmo3_instruct/upturn}{$-$0.09}
\statdef{olmo3_instruct/upturn_ci}{[$-$0.19, $+$0.01]}
\statdef{olmo3_instruct/lowest}{1.00}
\statdef{olmo3_instruct/end}{-0.5260}
\statdef{olmo3_instruct/stated_twin}{$-$0.83}
\statdef{olmo3_instruct/twin_implicit}{$+$0.40}
\statdef{olmo3_instruct/twin_template}{$+$0.13}
\statdef{olmo3_instruct/partial}{$+$0.05}
\statdef{olmo3_instruct/slope_w1}{$-$0.99}
\statdef{olmo3_instruct/twin_w1}{$+$0.08}
\statdef{olmo3_instruct/slope_w2}{$-$0.24}
\statdef{olmo3_instruct/twin_w2}{$+$0.09}
\statdef{olmo3_instruct/slope_w3}{$-$1.53}
\statdef{olmo3_instruct/twin_w3}{$+$0.09}
\statdef{olmo3_instruct/auc_stated_streaky}{0.77}
\statdef{olmo3_instruct/auc_implicit_streaky}{0.52}
\statdef{olmo3_instruct/auc_stated_alternating}{0.40}
\statdef{olmo3_instruct/auc_implicit_alternating}{0.57}
\statdef{olmo3_instruct/auc_stated_heads}{0.61}
\statdef{olmo3_instruct/auc_implicit_heads}{0.66}
\statdef{olmo3_think_sft/template}{$-$0.51}
\statdef{olmo3_think_sft/template_ci}{[$-$0.57, $-$0.45]}
\statdef{olmo3_think_sft/implicit_template}{$+$0.09}
\statdef{olmo3_think_sft/agree}{$-$0.49}
\statdef{olmo3_think_sft/agree_ci}{[$-$0.54, $-$0.43]}
\statdef{olmo3_think_sft/agree_hi}{$-$0.426}
\statdef{olmo3_think_sft/alt_fair}{0.441}
\statdef{olmo3_think_sft/alt_log}{0.450}
\statdef{olmo3_think_sft/coin}{$-$0.009}
\statdef{olmo3_think_sft/fair_excess}{$-$0.059}
\statdef{olmo3_think_sft/log_excess}{$-$0.050}
\statdef{olmo3_think_sft/mass}{0.75}
\statdef{olmo3_think_sft/vertex}{1.23}
\statdef{olmo3_think_sft/curv_lo}{0.7}
\statdef{olmo3_think_sft/ok}{3}
\statdef{olmo3_think_sft/upturn}{$-$0.96}
\statdef{olmo3_think_sft/upturn_ci}{[$-$1.17, $-$0.74]}
\statdef{olmo3_think_sft/lowest}{1.00}
\statdef{olmo3_think_sft/end}{-1.3819}
\statdef{olmo3_think_sft/stated_twin}{$-$0.71}
\statdef{olmo3_think_sft/twin_implicit}{$+$0.52}
\statdef{olmo3_think_sft/twin_template}{$+$0.06}
\statdef{olmo3_think_sft/partial}{$-$0.19}
\statdef{olmo3_think_sft/slope_w1}{$-$0.37}
\statdef{olmo3_think_sft/twin_w1}{$+$0.07}
\statdef{olmo3_think_sft/slope_w2}{$-$0.14}
\statdef{olmo3_think_sft/twin_w2}{$+$0.05}
\statdef{olmo3_think_sft/slope_w3}{$-$0.61}
\statdef{olmo3_think_sft/twin_w3}{$-$0.03}
\statdef{olmo3_think_sft/auc_stated_streaky}{0.76}
\statdef{olmo3_think_sft/auc_implicit_streaky}{0.52}
\statdef{olmo3_think_sft/auc_stated_alternating}{0.31}
\statdef{olmo3_think_sft/auc_implicit_alternating}{0.56}
\statdef{olmo3_think_sft/auc_stated_heads}{0.52}
\statdef{olmo3_think_sft/auc_implicit_heads}{0.65}
\statdef{olmo3_think_dpo/template}{$-$0.58}
\statdef{olmo3_think_dpo/template_ci}{[$-$0.64, $-$0.53]}
\statdef{olmo3_think_dpo/implicit_template}{$+$0.08}
\statdef{olmo3_think_dpo/agree}{$-$0.41}
\statdef{olmo3_think_dpo/agree_ci}{[$-$0.48, $-$0.35]}
\statdef{olmo3_think_dpo/agree_hi}{$-$0.348}
\statdef{olmo3_think_dpo/alt_fair}{0.443}
\statdef{olmo3_think_dpo/alt_log}{0.452}
\statdef{olmo3_think_dpo/coin}{$-$0.009}
\statdef{olmo3_think_dpo/fair_excess}{$-$0.057}
\statdef{olmo3_think_dpo/log_excess}{$-$0.048}
\statdef{olmo3_think_dpo/mass}{0.65}
\statdef{olmo3_think_dpo/vertex}{1.06}
\statdef{olmo3_think_dpo/curv_lo}{1.4}
\statdef{olmo3_think_dpo/ok}{3}
\statdef{olmo3_think_dpo/upturn}{$-$0.94}
\statdef{olmo3_think_dpo/upturn_ci}{[$-$1.12, $-$0.76]}
\statdef{olmo3_think_dpo/lowest}{1.00}
\statdef{olmo3_think_dpo/end}{-1.4026}
\statdef{olmo3_think_dpo/stated_twin}{$-$0.57}
\statdef{olmo3_think_dpo/twin_implicit}{$+$0.50}
\statdef{olmo3_think_dpo/twin_template}{$+$0.03}
\statdef{olmo3_think_dpo/partial}{$-$0.18}
\statdef{olmo3_think_dpo/slope_w1}{$-$0.45}
\statdef{olmo3_think_dpo/twin_w1}{$+$0.07}
\statdef{olmo3_think_dpo/slope_w2}{$-$0.13}
\statdef{olmo3_think_dpo/twin_w2}{$+$0.06}
\statdef{olmo3_think_dpo/slope_w3}{$-$0.71}
\statdef{olmo3_think_dpo/twin_w3}{$-$0.06}
\statdef{olmo3_think_dpo/auc_stated_streaky}{0.80}
\statdef{olmo3_think_dpo/auc_implicit_streaky}{0.53}
\statdef{olmo3_think_dpo/auc_stated_alternating}{0.30}
\statdef{olmo3_think_dpo/auc_implicit_alternating}{0.57}
\statdef{olmo3_think_dpo/auc_stated_heads}{0.59}
\statdef{olmo3_think_dpo/auc_implicit_heads}{0.66}
\statdef{olmo3_think/template}{$-$0.55}
\statdef{olmo3_think/template_ci}{[$-$0.60, $-$0.49]}
\statdef{olmo3_think/implicit_template}{$+$0.10}
\statdef{olmo3_think/agree}{$-$0.45}
\statdef{olmo3_think/agree_ci}{[$-$0.51, $-$0.39]}
\statdef{olmo3_think/agree_hi}{$-$0.386}
\statdef{olmo3_think/alt_fair}{0.444}
\statdef{olmo3_think/alt_log}{0.453}
\statdef{olmo3_think/coin}{$-$0.009}
\statdef{olmo3_think/fair_excess}{$-$0.056}
\statdef{olmo3_think/log_excess}{$-$0.047}
\statdef{olmo3_think/mass}{0.89}
\statdef{olmo3_think/vertex}{1.02}
\statdef{olmo3_think/curv_lo}{1.5}
\statdef{olmo3_think/ok}{3}
\statdef{olmo3_think/upturn}{$-$0.83}
\statdef{olmo3_think/upturn_ci}{[$-$1.02, $-$0.65]}
\statdef{olmo3_think/lowest}{1.00}
\statdef{olmo3_think/end}{-1.2993}
\statdef{olmo3_think/stated_twin}{$-$0.68}
\statdef{olmo3_think/twin_implicit}{$+$0.48}
\statdef{olmo3_think/twin_template}{$+$0.06}
\statdef{olmo3_think/partial}{$-$0.19}
\statdef{olmo3_think/slope_w1}{$-$0.65}
\statdef{olmo3_think/twin_w1}{$+$0.11}
\statdef{olmo3_think/slope_w2}{$-$0.24}
\statdef{olmo3_think/twin_w2}{$+$0.07}
\statdef{olmo3_think/slope_w3}{$-$1.01}
\statdef{olmo3_think/twin_w3}{$-$0.04}
\statdef{olmo3_think/auc_stated_streaky}{0.79}
\statdef{olmo3_think/auc_implicit_streaky}{0.52}
\statdef{olmo3_think/auc_stated_alternating}{0.31}
\statdef{olmo3_think/auc_implicit_alternating}{0.57}
\statdef{olmo3_think/auc_stated_heads}{0.56}
\statdef{olmo3_think/auc_implicit_heads}{0.67}
\statdef{h8/qwen_instruct/base}{$+$0.64}
\statdef{h8/qwen_instruct/post}{$+$0.21}
\statdef{h8/qwen_instruct/diff}{$-$0.428}
\statdef{h8/qwen_instruct/ci}{[$-$0.497, $-$0.357]}
\statdef{h8/qwen_instruct/hi}{$-$0.357}
\statdef{h8/llama_instruct/base}{$+$0.41}
\statdef{h8/llama_instruct/post}{$-$0.43}
\statdef{h8/llama_instruct/diff}{$-$0.839}
\statdef{h8/llama_instruct/ci}{[$-$0.938, $-$0.738]}
\statdef{h8/llama_instruct/hi}{$-$0.738}
\statdef{h8/olmo2_instruct/base}{$+$0.37}
\statdef{h8/olmo2_instruct/post}{$+$0.23}
\statdef{h8/olmo2_instruct/diff}{$-$0.139}
\statdef{h8/olmo2_instruct/ci}{[$-$0.233, $-$0.050]}
\statdef{h8/olmo2_instruct/hi}{$-$0.050}
\statdef{h8/olmo3_instruct/base}{$-$0.22}
\statdef{h8/olmo3_instruct/post}{$-$0.30}
\statdef{h8/olmo3_instruct/diff}{$-$0.085}
\statdef{h8/olmo3_instruct/ci}{[$-$0.169, $-$0.001]}
\statdef{h8/olmo3_instruct/hi}{$-$0.001}
\statdef{h8/olmo3_think/base}{$-$0.22}
\statdef{h8/olmo3_think/post}{$-$0.45}
\statdef{h8/olmo3_think/diff}{$-$0.229}
\statdef{h8/olmo3_think/ci}{[$-$0.312, $-$0.147]}
\statdef{h8/olmo3_think/hi}{$-$0.147}
\statdef{partial/qwen_instruct/base}{$+$0.61}
\statdef{partial/qwen_instruct/post}{$+$0.22}
\statdef{partial/qwen_instruct/diff}{$-$0.39}
\statdef{partial/qwen_instruct/ci}{[$-$0.46, $-$0.32]}
\statdef{partial/llama_instruct/base}{$+$0.18}
\statdef{partial/llama_instruct/post}{$-$0.47}
\statdef{partial/llama_instruct/diff}{$-$0.65}
\statdef{partial/llama_instruct/ci}{[$-$0.77, $-$0.52]}
\statdef{partial/olmo2_instruct/base}{$+$0.29}
\statdef{partial/olmo2_instruct/post}{$+$0.21}
\statdef{partial/olmo2_instruct/diff}{$-$0.08}
\statdef{partial/olmo2_instruct/ci}{[$-$0.18, $+$0.02]}
\statdef{partial/olmo3_instruct/base}{$+$0.26}
\statdef{partial/olmo3_instruct/post}{$+$0.05}
\statdef{partial/olmo3_instruct/diff}{$-$0.20}
\statdef{partial/olmo3_instruct/ci}{[$-$0.30, $-$0.10]}
\statdef{partial/olmo3_think/base}{$+$0.26}
\statdef{partial/olmo3_think/post}{$-$0.19}
\statdef{partial/olmo3_think/diff}{$-$0.44}
\statdef{partial/olmo3_think/ci}{[$-$0.54, $-$0.33]}
\statdef{partial/drops}{4}
\statdef{h6/olmo3_instruct/rise}{$+$0.002}
\statdef{h6/olmo3_instruct/dpo}{0.000}
\statdef{h6/olmo3_think/rise}{$-$0.029}
\statdef{h6/olmo3_think/dpo}{$-$0.030}
\statdef{h6/olmo2/rise}{$+$0.089}
\statdef{h8/drop_min}{0.08}
\statdef{h8/drop_max}{0.84}
\statdef{h8/margin}{0.001}
\statdef{h5/margin}{0.015}
\statdef{olmo2/max_slope}{0.04}
\statdef{mass/chat_min}{0.97}
\statdef{pilot/packed/face/split}{0.63}
\statdef{pilot/packed/face/null}{0.19}
\statdef{pilot/packed/zorp/split}{0.59}
\statdef{pilot/packed/zorp/null}{0.19}
\statdef{pilot/packed/X/split}{0.62}
\statdef{pilot/packed/X/null}{0.19}
\statdef{pilot/packed/cross}{0.98}
\statdef{pilot/spaced/face/split}{0.81}
\statdef{pilot/spaced/face/null}{0.19}
\statdef{pilot/spaced/zorp/split}{0.83}
\statdef{pilot/spaced/zorp/null}{0.18}
\statdef{pilot/spaced/X/split}{0.76}
\statdef{pilot/spaced/X/null}{0.19}
\statdef{pilot/spaced/cross}{0.94}

\statdef{diag/qwen_base/says_first_h}{$-$0.57}
\statdef{diag/qwen_base/twin_first_h}{$-$0.23}
\statdef{diag/qwen_base/exp_first_h}{$-$0.59}
\statdef{diag/qwen_base/says_h_frac}{$-$0.61}
\statdef{diag/qwen_base/twin_h_frac}{$+$0.11}
\statdef{diag/qwen_base/exp_h_frac}{$-$0.33}
\statdef{diag/qwen_base/says_imb}{$+$0.05}
\statdef{diag/qwen_base/exp_imb}{$+$0.22}
\statdef{diag/qwen_base/partial_first_h}{$+$0.46}
\statdef{diag/qwen_base/partial_ht}{$+$0.40}
\statdef{diag/qwen_base/partial_ht_ci}{[$+$0.33, $+$0.47]}
\statdef{diag/qwen_base/partial_shape}{$+$0.60}
\statdef{diag/qwen_base/partial_ht_twin}{$+$0.12}
\statdef{diag/qwen_base/stated_twin_ht}{$-$0.67}
\statdef{diag/qwen_base/twin_expects_ht}{$-$0.48}
\statdef{diag/qwen_instruct/says_first_h}{$-$0.20}
\statdef{diag/qwen_instruct/twin_first_h}{$-$0.70}
\statdef{diag/qwen_instruct/exp_first_h}{$+$0.16}
\statdef{diag/qwen_instruct/says_h_frac}{$-$0.30}
\statdef{diag/qwen_instruct/twin_h_frac}{$-$0.21}
\statdef{diag/qwen_instruct/exp_h_frac}{$+$0.03}
\statdef{diag/qwen_instruct/says_imb}{$+$0.32}
\statdef{diag/qwen_instruct/exp_imb}{$+$0.30}
\statdef{diag/qwen_instruct/partial_first_h}{$+$0.25}
\statdef{diag/qwen_instruct/partial_ht}{$+$0.26}
\statdef{diag/qwen_instruct/partial_ht_ci}{[$+$0.18, $+$0.33]}
\statdef{diag/qwen_instruct/partial_shape}{$-$0.10}
\statdef{diag/qwen_instruct/partial_ht_twin}{$+$0.26}
\statdef{diag/qwen_instruct/stated_twin_ht}{$+$0.05}
\statdef{diag/qwen_instruct/twin_expects_ht}{$+$0.13}
\statdef{diag/llama_base/says_first_h}{$-$0.61}
\statdef{diag/llama_base/twin_first_h}{$+$0.69}
\statdef{diag/llama_base/exp_first_h}{$-$0.85}
\statdef{diag/llama_base/says_h_frac}{$-$0.27}
\statdef{diag/llama_base/twin_h_frac}{$-$0.16}
\statdef{diag/llama_base/exp_h_frac}{$-$0.15}
\statdef{diag/llama_base/says_imb}{$-$0.27}
\statdef{diag/llama_base/exp_imb}{$+$0.19}
\statdef{diag/llama_base/partial_first_h}{$-$0.28}
\statdef{diag/llama_base/partial_ht}{$-$0.27}
\statdef{diag/llama_base/partial_ht_ci}{[$-$0.35, $-$0.18]}
\statdef{diag/llama_base/partial_shape}{$+$0.61}
\statdef{diag/llama_base/partial_ht_twin}{$-$0.28}
\statdef{diag/llama_base/stated_twin_ht}{$-$0.17}
\statdef{diag/llama_base/twin_expects_ht}{$-$0.02}
\statdef{diag/llama_instruct/says_first_h}{$-$0.88}
\statdef{diag/llama_instruct/twin_first_h}{$+$0.75}
\statdef{diag/llama_instruct/exp_first_h}{$+$0.46}
\statdef{diag/llama_instruct/says_h_frac}{$-$0.30}
\statdef{diag/llama_instruct/twin_h_frac}{$+$0.18}
\statdef{diag/llama_instruct/exp_h_frac}{$+$0.41}
\statdef{diag/llama_instruct/says_imb}{$-$0.05}
\statdef{diag/llama_instruct/exp_imb}{$+$0.24}
\statdef{diag/llama_instruct/partial_first_h}{$-$0.05}
\statdef{diag/llama_instruct/partial_ht}{$+$0.03}
\statdef{diag/llama_instruct/partial_ht_ci}{[$-$0.05, $+$0.10]}
\statdef{diag/llama_instruct/partial_shape}{$-$0.43}
\statdef{diag/llama_instruct/partial_ht_twin}{$-$0.07}
\statdef{diag/llama_instruct/stated_twin_ht}{$-$0.27}
\statdef{diag/llama_instruct/twin_expects_ht}{$-$0.35}
\statdef{diag/olmo2_base/says_first_h}{$-$0.61}
\statdef{diag/olmo2_base/twin_first_h}{$+$0.41}
\statdef{diag/olmo2_base/exp_first_h}{$-$0.81}
\statdef{diag/olmo2_base/says_h_frac}{$-$0.36}
\statdef{diag/olmo2_base/twin_h_frac}{$-$0.15}
\statdef{diag/olmo2_base/exp_h_frac}{$+$0.23}
\statdef{diag/olmo2_base/says_imb}{$-$0.11}
\statdef{diag/olmo2_base/exp_imb}{$+$0.02}
\statdef{diag/olmo2_base/partial_first_h}{$-$0.26}
\statdef{diag/olmo2_base/partial_ht}{$-$0.08}
\statdef{diag/olmo2_base/partial_ht_ci}{[$-$0.16, 0.00]}
\statdef{diag/olmo2_base/partial_shape}{$+$0.38}
\statdef{diag/olmo2_base/partial_ht_twin}{$-$0.04}
\statdef{diag/olmo2_base/stated_twin_ht}{$-$0.18}
\statdef{diag/olmo2_base/twin_expects_ht}{$+$0.20}
\statdef{diag/olmo2_sft/says_first_h}{$-$0.46}
\statdef{diag/olmo2_sft/twin_first_h}{$+$0.28}
\statdef{diag/olmo2_sft/exp_first_h}{$+$0.32}
\statdef{diag/olmo2_sft/says_h_frac}{$+$0.03}
\statdef{diag/olmo2_sft/twin_h_frac}{$+$0.70}
\statdef{diag/olmo2_sft/exp_h_frac}{$+$0.32}
\statdef{diag/olmo2_sft/says_imb}{$+$0.09}
\statdef{diag/olmo2_sft/exp_imb}{$+$0.09}
\statdef{diag/olmo2_sft/partial_first_h}{$+$0.50}
\statdef{diag/olmo2_sft/partial_ht}{$+$0.48}
\statdef{diag/olmo2_sft/partial_ht_ci}{[$+$0.42, $+$0.54]}
\statdef{diag/olmo2_sft/partial_shape}{$+$0.23}
\statdef{diag/olmo2_sft/partial_ht_twin}{$+$0.22}
\statdef{diag/olmo2_sft/stated_twin_ht}{$-$0.75}
\statdef{diag/olmo2_sft/twin_expects_ht}{$-$0.47}
\statdef{diag/olmo2_dpo/says_first_h}{$-$0.19}
\statdef{diag/olmo2_dpo/twin_first_h}{$+$0.21}
\statdef{diag/olmo2_dpo/exp_first_h}{$+$0.25}
\statdef{diag/olmo2_dpo/says_h_frac}{$+$0.14}
\statdef{diag/olmo2_dpo/twin_h_frac}{$+$0.57}
\statdef{diag/olmo2_dpo/exp_h_frac}{$+$0.34}
\statdef{diag/olmo2_dpo/says_imb}{$+$0.08}
\statdef{diag/olmo2_dpo/exp_imb}{$+$0.09}
\statdef{diag/olmo2_dpo/partial_first_h}{$+$0.37}
\statdef{diag/olmo2_dpo/partial_ht}{$+$0.33}
\statdef{diag/olmo2_dpo/partial_ht_ci}{[$+$0.26, $+$0.39]}
\statdef{diag/olmo2_dpo/partial_shape}{$+$0.27}
\statdef{diag/olmo2_dpo/partial_ht_twin}{$+$0.10}
\statdef{diag/olmo2_dpo/stated_twin_ht}{$-$0.73}
\statdef{diag/olmo2_dpo/twin_expects_ht}{$-$0.36}
\statdef{diag/olmo2_instruct/says_first_h}{$+$0.10}
\statdef{diag/olmo2_instruct/twin_first_h}{$+$0.19}
\statdef{diag/olmo2_instruct/exp_first_h}{$+$0.20}
\statdef{diag/olmo2_instruct/says_h_frac}{$+$0.28}
\statdef{diag/olmo2_instruct/twin_h_frac}{$+$0.47}
\statdef{diag/olmo2_instruct/exp_h_frac}{$+$0.32}
\statdef{diag/olmo2_instruct/says_imb}{$+$0.11}
\statdef{diag/olmo2_instruct/exp_imb}{$+$0.04}
\statdef{diag/olmo2_instruct/partial_first_h}{$+$0.22}
\statdef{diag/olmo2_instruct/partial_ht}{$+$0.15}
\statdef{diag/olmo2_instruct/partial_ht_ci}{[$+$0.08, $+$0.23]}
\statdef{diag/olmo2_instruct/partial_shape}{$+$0.21}
\statdef{diag/olmo2_instruct/partial_ht_twin}{$-$0.04}
\statdef{diag/olmo2_instruct/stated_twin_ht}{$-$0.58}
\statdef{diag/olmo2_instruct/twin_expects_ht}{$-$0.32}
\statdef{diag/olmo3_base/says_first_h}{$-$0.59}
\statdef{diag/olmo3_base/twin_first_h}{$+$0.79}
\statdef{diag/olmo3_base/exp_first_h}{$+$0.48}
\statdef{diag/olmo3_base/says_h_frac}{$-$0.65}
\statdef{diag/olmo3_base/twin_h_frac}{$+$0.30}
\statdef{diag/olmo3_base/exp_h_frac}{$+$0.19}
\statdef{diag/olmo3_base/says_imb}{$+$0.15}
\statdef{diag/olmo3_base/exp_imb}{$+$0.40}
\statdef{diag/olmo3_base/partial_first_h}{$+$0.09}
\statdef{diag/olmo3_base/partial_ht}{$+$0.20}
\statdef{diag/olmo3_base/partial_ht_ci}{[$+$0.13, $+$0.27]}
\statdef{diag/olmo3_base/partial_shape}{$-$0.41}
\statdef{diag/olmo3_base/partial_ht_twin}{$+$0.40}
\statdef{diag/olmo3_base/stated_twin_ht}{$-$0.57}
\statdef{diag/olmo3_base/twin_expects_ht}{$+$0.21}
\statdef{diag/olmo3_instruct_sft/says_first_h}{$-$0.60}
\statdef{diag/olmo3_instruct_sft/twin_first_h}{$+$0.80}
\statdef{diag/olmo3_instruct_sft/exp_first_h}{$+$0.58}
\statdef{diag/olmo3_instruct_sft/says_h_frac}{$-$0.40}
\statdef{diag/olmo3_instruct_sft/twin_h_frac}{$+$0.50}
\statdef{diag/olmo3_instruct_sft/exp_h_frac}{$+$0.12}
\statdef{diag/olmo3_instruct_sft/says_imb}{$+$0.28}
\statdef{diag/olmo3_instruct_sft/exp_imb}{$+$0.11}
\statdef{diag/olmo3_instruct_sft/partial_first_h}{$-$0.10}
\statdef{diag/olmo3_instruct_sft/partial_ht}{$-$0.11}
\statdef{diag/olmo3_instruct_sft/partial_ht_ci}{[$-$0.20, $-$0.03]}
\statdef{diag/olmo3_instruct_sft/partial_shape}{$-$0.42}
\statdef{diag/olmo3_instruct_sft/partial_ht_twin}{$-$0.18}
\statdef{diag/olmo3_instruct_sft/stated_twin_ht}{$-$0.83}
\statdef{diag/olmo3_instruct_sft/twin_expects_ht}{$+$0.02}
\statdef{diag/olmo3_instruct_dpo/says_first_h}{$-$0.65}
\statdef{diag/olmo3_instruct_dpo/twin_first_h}{$+$0.81}
\statdef{diag/olmo3_instruct_dpo/exp_first_h}{$+$0.57}
\statdef{diag/olmo3_instruct_dpo/says_h_frac}{$-$0.42}
\statdef{diag/olmo3_instruct_dpo/twin_h_frac}{$+$0.60}
\statdef{diag/olmo3_instruct_dpo/exp_h_frac}{0.00}
\statdef{diag/olmo3_instruct_dpo/says_imb}{$+$0.24}
\statdef{diag/olmo3_instruct_dpo/exp_imb}{$+$0.21}
\statdef{diag/olmo3_instruct_dpo/partial_first_h}{$+$0.12}
\statdef{diag/olmo3_instruct_dpo/partial_ht}{$+$0.06}
\statdef{diag/olmo3_instruct_dpo/partial_ht_ci}{[$-$0.02, $+$0.14]}
\statdef{diag/olmo3_instruct_dpo/partial_shape}{$-$0.37}
\statdef{diag/olmo3_instruct_dpo/partial_ht_twin}{$+$0.08}
\statdef{diag/olmo3_instruct_dpo/stated_twin_ht}{$-$0.54}
\statdef{diag/olmo3_instruct_dpo/twin_expects_ht}{0.00}
\statdef{diag/olmo3_instruct/says_first_h}{$-$0.67}
\statdef{diag/olmo3_instruct/twin_first_h}{$+$0.80}
\statdef{diag/olmo3_instruct/exp_first_h}{$+$0.58}
\statdef{diag/olmo3_instruct/says_h_frac}{$-$0.46}
\statdef{diag/olmo3_instruct/twin_h_frac}{$+$0.61}
\statdef{diag/olmo3_instruct/exp_h_frac}{$-$0.02}
\statdef{diag/olmo3_instruct/says_imb}{$+$0.23}
\statdef{diag/olmo3_instruct/exp_imb}{$+$0.23}
\statdef{diag/olmo3_instruct/partial_first_h}{$+$0.15}
\statdef{diag/olmo3_instruct/partial_ht}{$+$0.08}
\statdef{diag/olmo3_instruct/partial_ht_ci}{[$-$0.01, $+$0.16]}
\statdef{diag/olmo3_instruct/partial_shape}{$-$0.38}
\statdef{diag/olmo3_instruct/partial_ht_twin}{$+$0.10}
\statdef{diag/olmo3_instruct/stated_twin_ht}{$-$0.56}
\statdef{diag/olmo3_instruct/twin_expects_ht}{$+$0.01}
\statdef{diag/olmo3_think_sft/says_first_h}{$-$0.64}
\statdef{diag/olmo3_think_sft/twin_first_h}{$+$0.72}
\statdef{diag/olmo3_think_sft/exp_first_h}{$+$0.59}
\statdef{diag/olmo3_think_sft/says_h_frac}{$-$0.40}
\statdef{diag/olmo3_think_sft/twin_h_frac}{$+$0.48}
\statdef{diag/olmo3_think_sft/exp_h_frac}{$+$0.20}
\statdef{diag/olmo3_think_sft/says_imb}{$+$0.18}
\statdef{diag/olmo3_think_sft/exp_imb}{$+$0.19}
\statdef{diag/olmo3_think_sft/partial_first_h}{$-$0.17}
\statdef{diag/olmo3_think_sft/partial_ht}{$-$0.15}
\statdef{diag/olmo3_think_sft/partial_ht_ci}{[$-$0.23, $-$0.07]}
\statdef{diag/olmo3_think_sft/partial_shape}{$-$0.53}
\statdef{diag/olmo3_think_sft/partial_ht_twin}{$-$0.10}
\statdef{diag/olmo3_think_sft/stated_twin_ht}{$-$0.37}
\statdef{diag/olmo3_think_sft/twin_expects_ht}{$+$0.14}
\statdef{diag/olmo3_think_dpo/says_first_h}{$-$0.53}
\statdef{diag/olmo3_think_dpo/twin_first_h}{$+$0.71}
\statdef{diag/olmo3_think_dpo/exp_first_h}{$+$0.58}
\statdef{diag/olmo3_think_dpo/says_h_frac}{$-$0.36}
\statdef{diag/olmo3_think_dpo/twin_h_frac}{$+$0.50}
\statdef{diag/olmo3_think_dpo/exp_h_frac}{$+$0.20}
\statdef{diag/olmo3_think_dpo/says_imb}{$+$0.22}
\statdef{diag/olmo3_think_dpo/exp_imb}{$+$0.22}
\statdef{diag/olmo3_think_dpo/partial_first_h}{$-$0.15}
\statdef{diag/olmo3_think_dpo/partial_ht}{$-$0.13}
\statdef{diag/olmo3_think_dpo/partial_ht_ci}{[$-$0.22, $-$0.05]}
\statdef{diag/olmo3_think_dpo/partial_shape}{$-$0.48}
\statdef{diag/olmo3_think_dpo/partial_ht_twin}{$-$0.11}
\statdef{diag/olmo3_think_dpo/stated_twin_ht}{$-$0.21}
\statdef{diag/olmo3_think_dpo/twin_expects_ht}{$+$0.13}
\statdef{diag/olmo3_think/says_first_h}{$-$0.60}
\statdef{diag/olmo3_think/twin_first_h}{$+$0.68}
\statdef{diag/olmo3_think/exp_first_h}{$+$0.56}
\statdef{diag/olmo3_think/says_h_frac}{$-$0.39}
\statdef{diag/olmo3_think/twin_h_frac}{$+$0.53}
\statdef{diag/olmo3_think/exp_h_frac}{$+$0.22}
\statdef{diag/olmo3_think/says_imb}{$+$0.21}
\statdef{diag/olmo3_think/exp_imb}{$+$0.23}
\statdef{diag/olmo3_think/partial_first_h}{$-$0.16}
\statdef{diag/olmo3_think/partial_ht}{$-$0.14}
\statdef{diag/olmo3_think/partial_ht_ci}{[$-$0.22, $-$0.05]}
\statdef{diag/olmo3_think/partial_shape}{$-$0.51}
\statdef{diag/olmo3_think/partial_ht_twin}{$-$0.10}
\statdef{diag/olmo3_think/stated_twin_ht}{$-$0.35}
\statdef{diag/olmo3_think/twin_expects_ht}{$+$0.12}
\statdef{diag/first_h/qwen_instruct/base}{$+$0.46}
\statdef{diag/first_h/qwen_instruct/post}{$+$0.25}
\statdef{diag/first_h/qwen_instruct/diff}{$-$0.21}
\statdef{diag/first_h/qwen_instruct/ci}{[$-$0.28, $-$0.13]}
\statdef{diag/ht/qwen_instruct/base}{$+$0.40}
\statdef{diag/ht/qwen_instruct/post}{$+$0.26}
\statdef{diag/ht/qwen_instruct/diff}{$-$0.14}
\statdef{diag/ht/qwen_instruct/ci}{[$-$0.22, $-$0.06]}
\statdef{diag/ht_twin/qwen_instruct/base}{$+$0.12}
\statdef{diag/ht_twin/qwen_instruct/post}{$+$0.26}
\statdef{diag/ht_twin/qwen_instruct/diff}{$+$0.14}
\statdef{diag/ht_twin/qwen_instruct/ci}{[$+$0.04, $+$0.23]}
\statdef{diag/first_h/llama_instruct/base}{$-$0.28}
\statdef{diag/first_h/llama_instruct/post}{$-$0.05}
\statdef{diag/first_h/llama_instruct/diff}{$+$0.22}
\statdef{diag/first_h/llama_instruct/ci}{[$+$0.10, $+$0.35]}
\statdef{diag/ht/llama_instruct/base}{$-$0.27}
\statdef{diag/ht/llama_instruct/post}{$+$0.03}
\statdef{diag/ht/llama_instruct/diff}{$+$0.29}
\statdef{diag/ht/llama_instruct/ci}{[$+$0.17, $+$0.41]}
\statdef{diag/ht_twin/llama_instruct/base}{$-$0.28}
\statdef{diag/ht_twin/llama_instruct/post}{$-$0.07}
\statdef{diag/ht_twin/llama_instruct/diff}{$+$0.20}
\statdef{diag/ht_twin/llama_instruct/ci}{[$+$0.08, $+$0.31]}
\statdef{diag/first_h/olmo2_instruct/base}{$-$0.26}
\statdef{diag/first_h/olmo2_instruct/post}{$+$0.22}
\statdef{diag/first_h/olmo2_instruct/diff}{$+$0.48}
\statdef{diag/first_h/olmo2_instruct/ci}{[$+$0.37, $+$0.58]}
\statdef{diag/ht/olmo2_instruct/base}{$-$0.08}
\statdef{diag/ht/olmo2_instruct/post}{$+$0.15}
\statdef{diag/ht/olmo2_instruct/diff}{$+$0.23}
\statdef{diag/ht/olmo2_instruct/ci}{[$+$0.11, $+$0.36]}
\statdef{diag/ht_twin/olmo2_instruct/base}{$-$0.04}
\statdef{diag/ht_twin/olmo2_instruct/post}{$-$0.04}
\statdef{diag/ht_twin/olmo2_instruct/diff}{0.00}
\statdef{diag/ht_twin/olmo2_instruct/ci}{[$-$0.13, $+$0.13]}
\statdef{diag/first_h/olmo3_instruct/base}{$+$0.09}
\statdef{diag/first_h/olmo3_instruct/post}{$+$0.15}
\statdef{diag/first_h/olmo3_instruct/diff}{$+$0.06}
\statdef{diag/first_h/olmo3_instruct/ci}{[$-$0.04, $+$0.16]}
\statdef{diag/ht/olmo3_instruct/base}{$+$0.20}
\statdef{diag/ht/olmo3_instruct/post}{$+$0.08}
\statdef{diag/ht/olmo3_instruct/diff}{$-$0.13}
\statdef{diag/ht/olmo3_instruct/ci}{[$-$0.23, $-$0.03]}
\statdef{diag/ht_twin/olmo3_instruct/base}{$+$0.40}
\statdef{diag/ht_twin/olmo3_instruct/post}{$+$0.10}
\statdef{diag/ht_twin/olmo3_instruct/diff}{$-$0.31}
\statdef{diag/ht_twin/olmo3_instruct/ci}{[$-$0.40, $-$0.21]}
\statdef{diag/first_h/olmo3_think/base}{$+$0.09}
\statdef{diag/first_h/olmo3_think/post}{$-$0.16}
\statdef{diag/first_h/olmo3_think/diff}{$-$0.25}
\statdef{diag/first_h/olmo3_think/ci}{[$-$0.36, $-$0.15]}
\statdef{diag/ht/olmo3_think/base}{$+$0.20}
\statdef{diag/ht/olmo3_think/post}{$-$0.14}
\statdef{diag/ht/olmo3_think/diff}{$-$0.34}
\statdef{diag/ht/olmo3_think/ci}{[$-$0.44, $-$0.24]}
\statdef{diag/ht_twin/olmo3_think/base}{$+$0.40}
\statdef{diag/ht_twin/olmo3_think/post}{$-$0.10}
\statdef{diag/ht_twin/olmo3_think/diff}{$-$0.50}
\statdef{diag/ht_twin/olmo3_think/ci}{[$-$0.60, $-$0.40]}
\statdef{diag/ht/drops}{3}
\statdef{diag/first_h/drops}{2}

\makeatletter
\newcommand*{\rows}[1]{\@@input #1 }
\makeatother

\definecolor{sone}{HTML}{2A78D6}
\definecolor{stwo}{HTML}{EB6834}
\definecolor{sthree}{HTML}{1BAF7A}
\definecolor{inktwo}{HTML}{52514E}
\definecolor{muted}{HTML}{898781}
\definecolor{hair}{HTML}{E1E0D9}
\definecolor{rule}{HTML}{C3C2B7}
\definecolor{band}{HTML}{F0EFEC}

\pgfplotsset{
  paper/.style={
    axis lines*=left,
    axis line style={rule},
    tick style={rule},
    tick label style={font=\footnotesize, text=inktwo},
    label style={font=\footnotesize, text=inktwo},
    title style={font=\footnotesize\bfseries, text=black},
    major grid style={hair, line width=0.4pt},
    legend style={draw=none, fill=none, font=\footnotesize, text=inktwo},
    legend cell align=left,
  },
}

\captionsetup{font=small, labelfont=bf}
\setlist{itemsep=2pt, topsep=4pt}

\title{Saying Is Not Expecting: Post-Training Decouples What Language Models Say About Randomness from What They Expect}
\author{Zeyad Deeb\\ Independent researcher}
\date{November 2026}

\begin{document}
\maketitle

\begin{abstract}
People judge streaky coin sequences to be rigged and expect a fair coin to switch sides more often than it does. We put the same question to language models in two ways at once: what a model says when asked whether a 16-flip sequence came from a fair or a rigged coin, and what it expects, read from the likelihood it assigns the same flips under fair and rigged framings. In a pre-registered study of 15 open checkpoints from four families (Qwen2.5 3B, Llama 3.2 3B, OLMo 2 1B and OLMo 3 7B, including intermediate SFT and DPO checkpoints), the hypotheses, thresholds and analysis code were hashed before the held-out stimuli were generated. Every question has a twin with nonsense answer labels, and every expectation has a control that presents the same letters as a system log. Seven of eight registered hypotheses passed. In all five base-to-post-trained pairs, post-training lowered the coin-by-coin agreement between what the model says and what it expects, by \stat{h8/drop_min} to \stat{h8/drop_max} in correlation, in one pair by a margin of \stat{h8/margin}; with the nonsense twin partialled out, the drop holds in four of the five pairs. Much of the stated--twin correlation comes from a preference for heads over tails on the first flip; controlling for the first flip and heads share removes the drop in Llama and OLMo 2 (whose answers are not usable) and keeps it in Qwen and both OLMo 3 paths. Most of the extra switching that post-trained models expect from a fair coin also appears for the system-log framing, so it is not a belief about coins. Our prediction that OLMo 3 acquires this repetition avoidance during preference tuning failed: it never does, and the OLMo 3 reasoning checkpoints expect streaks while saying that streaks look rigged. A pilot in the style of superstitious-perception experiments shows why the nonsense twin is necessary: with exact log-probabilities, reverse-correlation maps are just as reproducible for a nonsense question as for a real one.
\end{abstract}

\section{Introduction}

Asked which sequences came from a fair coin, people favor sequences that switch more often than chance. Streaks look rigged, and the sequences people rate most random switch about 60\% of the time rather than 50\% \citep{falk1997making,barhillel1991perception}. The same intuition underlies the gambler's fallacy and the belief in a hot hand \citep{kahneman1972subjective,gilovich1985hot}, and it has rational readings: in short windows of a fair process, streaks really are rarer than people's long-run intuitions suggest \citep{hahn2009perceptions}. \citet{griffiths2003algorithmic} cast subjective randomness as a likelihood ratio: a sequence looks random to the extent that a random process explains it better than a regular one.

A language model offers two separate windows onto the same question. We can ask it, in words, whether a sequence came from a fair or a rigged coin. We can also read its expectation directly: the probability it assigns the sequence when told the coin is fair, against the probability when told the coin is rigged, is the model's own version of that likelihood ratio. If a model's verbal judgment reports its expectations, the two should agree coin by coin. Work on metalinguistic prompting \citep{hu2023prompting} and on generator--validator consistency \citep{li2024benchmarking,west2023generative} gives reasons to doubt it, but it has not traced where in training the two come apart, and randomness is an unusually clean place to look, because the stimuli are pure noise and the correct answer is known.

We measure both, on the same sequences, for 15 checkpoints. They include the intermediate SFT and DPO checkpoints of OLMo 2 1B and both post-training paths of OLMo 3 7B \citep{olmo2024olmo2,olmo2025olmo3}, so changes can be located in training. Two controls make the measures interpretable. Every rigged-or-fair question has a twin that asks the same thing with nonsense labels (``zorp'' or ``blarg''), which exposes any dependence on the sequence that has nothing to do with randomness. Every expectation measured under a coin framing is also measured for the same letters presented as a system log, which separates beliefs about coins from a generic reluctance to repeat a token. Hypotheses, thresholds and analysis code were sealed with a hash before the evaluation stimuli existed (\S\ref{sec:prereg}).

Our findings:
\begin{itemize}
  \item \textbf{Post-training decouples saying from expecting.} Agreement between stated judgment and likelihood ratio fell in all five pairs (H8). Controlling for which side lands first and the heads share removes the drop in Llama and OLMo 2, whose answers are not usable, and leaves it in Qwen and both OLMo 3 paths (\S\ref{sec:agreement}).
  \item \textbf{Over-alternation is mostly generic.} Post-trained models expect a fair coin to switch more than half the time, but nearly as much when the same letters are framed as a log file. The coin-specific part is at most \stat{olmo2_instruct/coin} in any post-trained checkpoint (H4, \S\ref{sec:switching}).
  \item \textbf{A prediction failed.} From OLMo 2 we predicted that OLMo 3 would acquire this repetition avoidance by DPO. It never acquires it, and its reasoning checkpoints expect streaks while saying streaks look rigged (H6).
  \item \textbf{Reproducible is not real.} Exact log-probabilities make any probe reproducible, including nonsense ones (\S\ref{sec:trap}). A stated measure is only interpretable beside a twin that cannot carry the concept.
\end{itemize}

\section{Related work}

\paragraph{Randomness in language models.} \citet{bigelow2024incontext} used random binary sequences to study in-context learning in GPT-3.5-class models and found sharp transitions from seemingly random output to deterministic repetition. \citet{vankoevering2024heads} found that coin flips generated by language models reproduce human biases, including over-alternation and aversion to long runs, often in amplified form, and that over-alternation is not explained by memorization or tokenization alone. Both study generation. We study judgment and likelihood, and we add a framing of the same letters that is not about coins.

\paragraph{Saying versus knowing.} \citet{hu2023prompting} found that metalinguistic judgments elicited by prompts are worse than quantities read directly from a model's probabilities, increasingly so as the prompt moves away from next-word prediction, and \citet{hu2024auxiliary} showed that such task demands mask abilities most in small models. Generator--validator gaps \citep{li2024benchmarking,west2023generative} and self-evaluation \citep{kadavath2022language} concern whether a model's answers about its outputs match the outputs. \citet{cacioli2026signal} applied signal detection theory to factual discrimination in a pre-registered study and found more asymmetric evidence distributions in instruct models than in a base model. We compare a model's stated judgment with its own likelihoods on the same items and follow the comparison through post-training stages.

\paragraph{Reverse correlation and multiple choice.} Superstitious-perception experiments show observers pure noise, ask them to detect a target, and average the noise by response to reveal an internal template \citep{gosselin2003superstitious,liu2014seeing,murray2011classification}. Our pilot transferred the method to language models and found its main pitfall (\S\ref{sec:trap}). Language models prefer some answer positions and letters regardless of content \citep{zheng2024large}, so every stated score here is counterbalanced over the order of the options.

\section{Method}

\subsection{Stimuli}
Each stimulus is a sequence of 16 flips written as space-separated letters (\texttt{H T T H \ldots}). Evaluation seed \stat{seed_eval} produced \stat{n_fair} fair sequences from a SplitMix64 generator; \stat{n_strata} stratified sequences with exactly $k$ switches for each $k = 0, \ldots, 15$ (\stat{n_per_k} per $k$, switch positions uniform); and three rigged sets of \stat{n_rigged}: \emph{streaky} (a Markov chain that repeats the last outcome with probability 0.8), \emph{alternating} (repeats with probability 0.2) and \emph{biased} (heads with probability 0.85). The switch rate of a sequence is $k/15$.

\subsection{What a model says}
Three wordings (Appendix~\ref{app:prompts}) show the sequence, say that fair and rigged sources are equally common, and offer two options labeled A and B, followed by ``Reply with a single letter (A or B).'' Post-trained checkpoints receive the prompt through their own chat template with the system message ``Answer with only A or B.''; base checkpoints receive the prompt followed by a newline and ``Answer:''. The OLMo 3 Think checkpoints open a thought in their template, which we close empty with \texttt{</think>} and a blank line before reading the answer; that choice was fixed on calibration data by the probability mass it leaves on A and B.

With $\ell = \log P(\mathrm{A}) - \log P(\mathrm{B})$, where each letter's probability sums its single-token forms with and without a leading space, the stated score of sequence $x$ under wording $w$ is
\[
s(x, w) = \tfrac12 \left[ \ell(\text{rigged listed as A}) - \ell(\text{rigged listed as B}) \right],
\]
so $s > 0$ leans rigged and any constant preference for a letter cancels. The nonsense twin keeps the lead and the sequence but asks whether the sequence is ``a zorp'' or ``a blarg''; its score $z(x, w)$ is built the same way. The labels mean nothing, so any dependence of $z$ on the sequence reflects how the model reads sequences through this format, not a concept of randomness. A checkpoint's stated judgment is \emph{usable} when the least-squares slope of $s$ on switch rate over $k = 1, \ldots, 14$ is at least 0.10 logits per unit switch rate and at least twice the twin's slope, in at least two of the three wordings.

\subsection{What a model expects}
Expectations are read from raw-text continuations for every checkpoint, without chat templates. For a framing pair $f$ (for example ``Here are 16 flips of a fair coin'' and ``\ldots of a rigged coin''), the implicit score is the likelihood ratio
\[
i(x, f) = \log p(x \mid \text{rigged framing}) - \log p(x \mid \text{fair framing}),
\]
summed over the 16 flip tokens; $i > 0$ means the sequence is more expected from a rigged coin. Expected switching $a(x, c)$ is the mean, over flips 2 to 16, of the probability the model gives to switching, renormalized over the two letters, under context $c$. A model that treats the coin as fair scores exactly 0.5; because $a$ is a predictive probability rather than a frequency of selected flips, it does not inherit the streak-selection bias of finite sequences \citep{miller2018surprised}. We compute $a$ under the three fair framings and under two framings that present the same letters as a system log or as machine status codes. The coin-specific excess is the mean of $a$ under the fair framings minus its mean under the log framings.

\subsection{Summary statistics}
Pooled stated and implicit scores average each coin's scores over wordings or framings after standardizing within each. A \emph{template} is the Pearson correlation of a pooled score with switch rate across the \stat{n_fair} fair coins; a negative stated template means streaky sequences are called rigged. \emph{Agreement} is the correlation between the pooled stated and pooled implicit scores across fair coins. The stated \emph{curve} is the mean pooled stated score at each $k$; it was registered as human-like when a quadratic fit has positive curvature (95\% interval above 0) and a vertex at a switch rate in $[0.55, 0.70]$. Intervals are percentile intervals from \stat{boot} bootstrap resamples of coins, paired across checkpoints for differences.

\subsection{Models and computation}
We score 15 checkpoints at pinned revisions (Appendix~\ref{app:models}): Qwen2.5 3B and its Instruct model \citep{qwen2024qwen25}; Llama 3.2 3B and its Instruct model \citep{grattafiori2024llama}; OLMo 2 1B base, SFT, DPO and the released Instruct model trained further with RL from verifiable rewards \citep{olmo2024olmo2,lambert2024tulu,rafailov2023direct}; and OLMo 3 7B base with the SFT, DPO and released final checkpoints of both its Instruct and its Think path \citep{olmo2025olmo3}. All scoring ran in 32-bit floating point on an Apple-silicon GPU with Transformers 5.23 and PyTorch 2.15. We rejected bfloat16 because it rounds the log-odds to steps of 0.0625. The prompts for one coin share a prefix, which is computed once and reused; on this setup the reuse is bit-identical to scoring each prompt in full.

\subsection{Pre-registration}
\label{sec:prereg}
Hypotheses were written from calibration screens on seed 0, which uses different stimuli, and then sealed \citep{nosek2018preregistration}. The protocol file holds the hypotheses, thresholds, wordings, model revisions and the SHA-256 of every analysis file; its own hash, \texttt{\stat{sha12}}, names the run, and the runner refuses the evaluation seed unless the protocol is marked sealed. Seed \stat{seed_eval} had never been generated before the sealed run.

Calibration changed hypotheses before sealing, and the protocol records how. H1 and H2 originally covered both 3B instruct models; exact scoring showed that Llama's stated answers were not usable, so they were narrowed to Qwen. H7 originally predicted that no model's curve would be human-like; Qwen contradicted it on calibration data, and it was reversed. H8 was added after calibration showed the agreement drop in the Qwen, Llama and OLMo 2 pairs, so for those three pairs H8 is a held-out replication, and only the two OLMo 3 pairs were predictions. H6 and all OLMo 3 results are predictions; no OLMo 3 checkpoint was calibrated. After the evaluation we recorded three amendments in a separate file, because the protocol cannot be edited without changing its hash. None changes a verdict (Appendix~\ref{app:amendments}).

\section{Results}

\begin{table}[t]
\centering
\small
\caption{Registered hypotheses, scored on held-out seed \stat{seed_eval} exactly as written.}
\label{tab:verdicts}
\begin{tabularx}{\textwidth}{@{}l>{\raggedright\arraybackslash}X>{\raggedright\arraybackslash}X>{\raggedright\arraybackslash}Xl@{}}
\toprule
 & Claim & Measured & Registered bar & Result \\
\midrule
\rows{generated/verdicts}
\bottomrule
\end{tabularx}
\end{table}

Seven of the eight registered hypotheses passed and one failed (Table~\ref{tab:verdicts}). Three passes are close to their bars: H4 met its threshold exactly, H5 cleared its ceiling by \stat{h5/margin}, and one of the five H8 pairs cleared zero by \stat{h8/margin}. Table~\ref{tab:checkpoints} gives the main measures for every checkpoint.

\begin{table}[t]
\centering
\footnotesize
\caption{Per-checkpoint measures. Usable: wordings in which the stated judgment passes the usability rule. Template: correlation of the pooled stated score with switch rate across fair coins (negative: streaks look rigged). Agreement: correlation between stated and implicit scores across fair coins, with 95\% interval. Expected switching under coin and log framings, and their difference. A/B mass: median probability on A or B at the answer position.}
\label{tab:checkpoints}
\setlength{\tabcolsep}{4pt}
\begin{tabular}{@{}lccccccc@{}}
\toprule
 & & & & \multicolumn{2}{c}{Expected switching} & & \\
\cmidrule(lr){5-6}
Checkpoint & Usable & Template & Agreement & Coin & Log & \makecell{Coin-\\specific} & \makecell{A/B\\mass} \\
\midrule
\rows{generated/checkpoints}
\bottomrule
\end{tabular}
\end{table}

\subsection{Saying and expecting come apart}
\label{sec:agreement}

\begin{figure}[t]
\centering
\pgfplotstableread[col sep=comma]{generated/pairs.csv}\pairs
\pgfplotstableread[col sep=comma]{generated/partial.csv}\partials
\begin{tikzpicture}
\begin{groupplot}[
  group style={group size=2 by 1, horizontal sep=0.6cm, y descriptions at=edge left},
  paper, scale only axis, width=5.4cm, height=3.9cm,
  xmin=-0.7, xmax=0.8, ymin=-0.6, ymax=4.6,
  ytick={0,1,2,3,4}, yticklabels from table={\pairs}{label},
  ytick style={draw=none}, y axis line style={draw=none},
  xtick={-0.5,0,0.5}, xticklabels={$-0.5$, $0$, $+0.5$}, xmajorgrids,
]
\nextgroupplot[title={As registered (H8)}, xlabel={$r$(stated, implicit) across fair coins}, legend to name=pairlegend, legend columns=2, legend style={/tikz/every even column/.append style={column sep=0.6cm}}]
\addplot[rule, densely dashed, forget plot] coordinates {(0,-0.6) (0,4.6)};
\addplot[inktwo, line width=0.9pt, quiver={u=\thisrow{diff}, v=0}, forget plot] table[x=base, y=idx]{\pairs};
\addplot[only marks, mark=*, mark size=2.6pt, color=sone, mark options={draw=white, line width=0.5pt}] table[x=base, y=idx]{\pairs};
\addlegendentry{base model}
\addplot[only marks, mark=square*, mark size=2.3pt, color=stwo, mark options={draw=white, line width=0.5pt}] table[x=post, y=idx]{\pairs};
\addlegendentry{post-trained}
\nextgroupplot[title={Nonsense twin partialled out (post hoc)}, xlabel={partial $r$, controlling for the twin}]
\addplot[rule, densely dashed, forget plot] coordinates {(0,-0.6) (0,4.6)};
\addplot[inktwo, line width=0.9pt, quiver={u=\thisrow{diff}, v=0}, forget plot] table[x=base, y=idx]{\partials};
\addplot[only marks, mark=*, mark size=2.6pt, color=sone, mark options={draw=white, line width=0.5pt}] table[x=base, y=idx]{\partials};
\addplot[only marks, mark=square*, mark size=2.3pt, color=stwo, mark options={draw=white, line width=0.5pt}] table[x=post, y=idx]{\partials};
\end{groupplot}
\end{tikzpicture}

\vspace{2pt}
\ref{pairlegend}
\caption{Agreement between what a checkpoint says about each fair coin and what it expects from it, before (circle) and after (square) post-training. Left: the registered measure; all five differences have 95\% intervals below zero (Table~\ref{tab:pairs}). Right: the same correlation with the nonsense twin partialled out of both scores; the drop holds in four pairs and not for OLMo 2 1B.}
\label{fig:pairs}
\end{figure}

\begin{table}[t]
\centering
\small
\caption{Agreement between stated and implicit scores across fair coins, base and post-trained, with the paired bootstrap difference.}
\label{tab:pairs}
\begin{tabular}{@{}lcccc@{}}
\toprule
Pair & Base & Post-trained & Difference [95\% CI] & With twin partialled out \\
\midrule
Qwen2.5 3B & \stat{h8/qwen_instruct/base} & \stat{h8/qwen_instruct/post} & \stat{h8/qwen_instruct/diff} \stat{h8/qwen_instruct/ci} & \stat{partial/qwen_instruct/diff} \stat{partial/qwen_instruct/ci} \\
Llama 3.2 3B & \stat{h8/llama_instruct/base} & \stat{h8/llama_instruct/post} & \stat{h8/llama_instruct/diff} \stat{h8/llama_instruct/ci} & \stat{partial/llama_instruct/diff} \stat{partial/llama_instruct/ci} \\
OLMo 2 1B & \stat{h8/olmo2_instruct/base} & \stat{h8/olmo2_instruct/post} & \stat{h8/olmo2_instruct/diff} \stat{h8/olmo2_instruct/ci} & \stat{partial/olmo2_instruct/diff} \stat{partial/olmo2_instruct/ci} \\
OLMo 3 7B Instruct & \stat{h8/olmo3_instruct/base} & \stat{h8/olmo3_instruct/post} & \stat{h8/olmo3_instruct/diff} \stat{h8/olmo3_instruct/ci} & \stat{partial/olmo3_instruct/diff} \stat{partial/olmo3_instruct/ci} \\
OLMo 3 7B Think & \stat{h8/olmo3_think/base} & \stat{h8/olmo3_think/post} & \stat{h8/olmo3_think/diff} \stat{h8/olmo3_think/ci} & \stat{partial/olmo3_think/diff} \stat{partial/olmo3_think/ci} \\
\bottomrule
\end{tabular}
\end{table}

In three base models, what the model says about a coin goes with what it expects from it: agreement is \stat{h8/qwen_instruct/base} for Qwen, \stat{h8/llama_instruct/base} for Llama and \stat{h8/olmo2_instruct/base} for OLMo 2. OLMo 3 base starts negative, at \stat{h8/olmo3_instruct/base}. Post-training lowered agreement in every pair (Figure~\ref{fig:pairs}, Table~\ref{tab:pairs}), most for Llama, whose Instruct model says the opposite of what it expects (\stat{h8/llama_instruct/post}), and least for OLMo 3 Instruct, whose interval ends at \stat{h8/olmo3_instruct/hi}. H8 passed with five of five pairs, one of them by a margin no larger than the resampling noise in a bootstrap bound. H5 also passed: Qwen Instruct's agreement is \stat{qwen_instruct/agree}, with an upper bound of \stat{qwen_instruct/agree_hi} against the registered ceiling of 0.30, and Llama Instruct's is negative. Llama Instruct's stated judgment is not usable (\S\ref{sec:usable}), so its negative agreement means only that the weak signal in its answers runs against its likelihoods.

The stated score and the nonsense twin are not independent. Across fair coins they correlate at between \stat{olmo3_think_dpo/stated_twin} and \stat{olmo3_instruct_sft/stated_twin} for every OLMo 3 checkpoint and at \stat{llama_instruct/stated_twin} for Llama Instruct (Appendix Table~\ref{tab:readout}), so much of the coin-to-coin variation in what these models say is shared with a question that has no meaning. To check that the agreement drop is not carried by that shared part, we partialled the twin out of both scores after the fact. The drop remains in four of five pairs (Figure~\ref{fig:pairs}, right). It does not survive for OLMo 2 1B (\stat{partial/olmo2_instruct/diff} \stat{partial/olmo2_instruct/ci}), none of whose checkpoints gives a usable answer. With the twin removed, OLMo 3 base agrees positively (\stat{partial/olmo3_instruct/base}), and the marginal OLMo 3 Instruct pair becomes a clear drop (\stat{partial/olmo3_instruct/diff} \stat{partial/olmo3_instruct/ci}). If the twin carries a faint copy of the randomness judgment rather than pure readout noise, partialling it out removes some real signal too, so this check is conservative.

\paragraph{First-flip bias and usability (post hoc).} Three of the five base models (Qwen, Llama and OLMo 2) and two post-trained models (Llama and OLMo 2) are not usable (\S\ref{sec:usable}; slopes $\leq 0.04$ logits). Because pooled scores are standardized within each wording, positive agreement in those base models comes from a faint raw-text residual, not a verbal judgment. Checking which features of the coin drive the scores also explains the strong twin correlations (Appendix Table~\ref{tab:features}): whether flip~1 is \texttt{H} or \texttt{T} correlates at \stat{diag/llama_instruct/says_first_h} with Llama Instruct's stated score, \stat{diag/olmo3_instruct/says_first_h} with OLMo 3 Instruct's, \stat{diag/olmo3_instruct/twin_first_h} with its twin, and \stat{diag/llama_base/exp_first_h} and \stat{diag/olmo2_base/exp_first_h} with the implicit scores of Llama base and OLMo 2 base. A fair coin is symmetric under swapping \texttt{H} and \texttt{T}, so any dependence on flip~1 or heads share is bias, not a view of randomness.

Partialling out flip~1 and heads share flips both unusable pairs (Appendix Table~\ref{tab:symmetry}): Llama moves from \stat{diag/ht/llama_instruct/base} to \stat{diag/ht/llama_instruct/post} (\stat{diag/ht/llama_instruct/diff} \stat{diag/ht/llama_instruct/ci}) and OLMo 2 from \stat{diag/ht/olmo2_instruct/base} to \stat{diag/ht/olmo2_instruct/post} (\stat{diag/ht/olmo2_instruct/diff} \stat{diag/ht/olmo2_instruct/ci}). All three pairs with a usable post-trained model still fall: Qwen by \stat{diag/ht/qwen_instruct/diff} \stat{diag/ht/qwen_instruct/ci}, OLMo 3 Instruct by \stat{diag/ht/olmo3_instruct/diff} \stat{diag/ht/olmo3_instruct/ci}, and OLMo 3 Think by \stat{diag/ht/olmo3_think/diff} \stat{diag/ht/olmo3_think/ci}. Partialling out the twin as well leaves the drop only in the two OLMo 3 pairs (\stat{diag/ht_twin/olmo3_instruct/diff} and \stat{diag/ht_twin/olmo3_think/diff}), the only family usable at both ends. In Qwen Instruct, partialling out switch rate and heads--tails imbalance removes its remaining agreement (\stat{diag/qwen_instruct/partial_shape}).

\subsection{Which checkpoints can answer}
\label{sec:usable}
Qwen Instruct's stated judgment is usable in all three wordings, with slopes of \stat{qwen_instruct/slope_w1}, \stat{qwen_instruct/slope_w2} and \stat{qwen_instruct/slope_w3} logits per unit switch rate against twin slopes of \stat{qwen_instruct/twin_w1}, \stat{qwen_instruct/twin_w2} and \stat{qwen_instruct/twin_w3} (H1; Appendix Table~\ref{tab:slopes}). Neither Qwen base nor Llama base is usable in any wording (H3). Llama Instruct is usable in none either: its largest slope, \stat{llama_instruct/slope_w1}, is less than twice its twin's \stat{llama_instruct/twin_w1}. No OLMo 2 1B checkpoint is usable; their slopes never exceed \stat{olmo2/max_slope} logits in magnitude, in line with the larger cost of task demands for small models \citep{hu2024auxiliary}. Every OLMo 3 checkpoint is usable, including the base model in two of three wordings, the only base model to manage it. OLMo 3's base checkpoint has been through a mid-training stage \citep{olmo2025olmo3}, so ``base'' here does not mean unexposed to question-and-answer formats; we did not investigate further. A/B mass is at least \stat{mass/chat_min} for every chat checkpoint but lower for the Think checkpoints (\stat{olmo3_think_dpo/mass} to \stat{olmo3_think/mass}), whose stated scores are therefore conditional on an answer they put less weight on.

\subsection{What looks rigged}
\label{sec:curves}

\begin{figure}[t]
\centering
\begin{tikzpicture}
\begin{axis}[
  paper, scale only axis, width=8.6cm, height=5cm,
  xmin=0, xmax=1, ymin=-1.5, ymax=3.4,
  xlabel={switch rate of the sequence ($k/15$)},
  ylabel={stated lean toward ``rigged'' (pooled SD)},
  xtick={0,0.2,0.4,0.6,0.8,1}, ytick={-1,0,1,2,3},
  xticklabel style={/pgf/number format/fixed},
  ymajorgrids, clip=false,
  legend style={at={(0.5,-0.2)}, anchor=north, legend columns=3, /tikz/every even column/.append style={column sep=0.5cm}},
]
\addplot[fill=band, draw=none, forget plot] coordinates {(0.55,-1.5) (0.70,-1.5) (0.70,3.4) (0.55,3.4)} --cycle;
\node[font=\scriptsize, text=inktwo, anchor=south, align=center] at (axis cs:0.625,3.4) {registered\\human range};
\addplot[muted, densely dotted, forget plot] coordinates {(0,0) (1,0)};
\addplot[sone, line width=1pt, mark=*, mark size=1.6pt] table[col sep=comma, x=rate, y=qweninstruct]{generated/curves.csv};
\addlegendentry{Qwen2.5 3B Instruct}
\addplot[stwo, line width=1pt, mark=square*, mark size=1.5pt] table[col sep=comma, x=rate, y=olmo3instruct]{generated/curves.csv};
\addlegendentry{OLMo 3 7B Instruct}
\addplot[sthree, line width=1pt, mark=triangle*, mark size=2pt] table[col sep=comma, x=rate, y=olmo3think]{generated/curves.csv};
\addlegendentry{OLMo 3 7B Think}
\node[font=\scriptsize, text=inktwo, anchor=west] at (axis cs:1.02,\stat{qwen_instruct/end}) {Qwen Instruct};
\node[font=\scriptsize, text=inktwo, anchor=west] at (axis cs:1.02,\stat{olmo3_instruct/end}) {OLMo 3 Instruct};
\node[font=\scriptsize, text=inktwo, anchor=west] at (axis cs:1.02,\stat{olmo3_think/end}) {OLMo 3 Think};
\end{axis}
\end{tikzpicture}
\caption{Mean stated score by number of switches, for the three usable post-trained checkpoints at the end of their training paths. Higher means the sequence is called rigged more readily. All three call all-heads and all-tails sequences rigged. Only Qwen Instruct turns back up for perfect alternation; OLMo 3 Instruct flattens out, and OLMo 3 Think finds perfect alternation the most random of all. Values are in Appendix Table~\ref{tab:curves}.}
\label{fig:curves}
\end{figure}

Qwen Instruct says streaky coins are rigged: its pooled stated score correlates with switch rate at \stat{qwen_instruct/template} \stat{qwen_instruct/template_ci} (H2). Its expectations lean the same way (implicit template \stat{qwen_instruct/implicit_template}), yet coin by coin the two agree only weakly, so they share little beyond switch rate. The OLMo 3 checkpoints say the same thing, with stated templates from \stat{olmo3_instruct/template} to \stat{olmo3_think_dpo/template}, while after post-training their expectations barely depend on switch rate (implicit templates from \stat{olmo3_instruct/implicit_template} to \stat{olmo3_instruct_sft/implicit_template}).

The registered curve test passed for Qwen Instruct: a quadratic fit has positive curvature and a vertex at switch rate \stat{qwen_instruct/vertex}, inside the registered human range (H7). The full curve is less human than that summary (Figure~\ref{fig:curves}). All-heads and all-tails sequences look overwhelmingly rigged, the score falls to a broad floor at middle-to-high switch rates, with its lowest point at \stat{qwen_instruct/lowest}, and it rises again for perfect alternation. That last rise is real: in a post-hoc check, perfect alternation scores \stat{qwen_instruct/upturn} \stat{qwen_instruct/upturn_ci} pooled standard units higher than switch rates of 0.60 and 0.67. Qwen Instruct shares the human view that perfect alternation is suspicious, but the switch rate it finds most random lies above the 0.6 that people prefer. The registered test also turned out to be lenient: it accepts OLMo 3 Instruct DPO and final, whose curves fall and flatten without rising again (upturn \stat{olmo3_instruct/upturn} \stat{olmo3_instruct/upturn_ci} for the final checkpoint). The Think checkpoints find perfect alternation the most random of all (\stat{olmo3_think/upturn} \stat{olmo3_think/upturn_ci}).

Stated judgments catch streaky coins: the area under the ROC curve for streaky against fair coins is \stat{qwen_instruct/auc_stated_streaky} for Qwen Instruct and between \stat{olmo3_think_sft/auc_stated_streaky} and \stat{olmo3_think_dpo/auc_stated_streaky} for the post-trained OLMo 3 checkpoints (exploratory; Appendix Table~\ref{tab:sensitivity}). The same checkpoints rate alternating rigged coins as fairer than fair ones (\stat{qwen_instruct/auc_stated_alternating} for Qwen Instruct, \stat{olmo3_think/auc_stated_alternating} for OLMo 3 Think), as the human preference for alternation would predict. Expectations separate streaky from fair coins less well, with implicit AUCs between \stat{olmo3_instruct_sft/auc_implicit_streaky} and \stat{llama_base/auc_implicit_streaky}.

\subsection{Expected switching}
\label{sec:switching}

\begin{figure}[t]
\centering
\begin{tikzpicture}
\begin{groupplot}[
  group style={group size=5 by 1, horizontal sep=0.32cm, y descriptions at=edge left},
  paper, scale only axis, height=3.5cm,
  ymin=0.42, ymax=0.64, ytick={0.45,0.5,0.55,0.6},
  yticklabel style={/pgf/number format/fixed, /pgf/number format/precision=2, /pgf/number format/fixed zerofill},
  ymajorgrids, ylabel={expected $P(\text{switch})$},
  title style={align=center},
  xticklabel style={font=\scriptsize, rotate=40, anchor=north east, inner sep=1pt},
]
\nextgroupplot[title={Qwen2.5\\3B}, width=1.5cm, xmin=-0.5, xmax=1.5, xtick={0,1}, xticklabels={base,Inst.}, legend to name=altlegend, legend columns=3, legend style={/tikz/every even column/.append style={column sep=0.6cm}}]
\addplot[muted, densely dotted, forget plot] coordinates {(-0.5,0.5) (1.5,0.5)};
\addplot[sone, line width=1pt, mark=*, mark size=1.7pt] table[col sep=comma, x=idx, y=fair]{generated/stages-qwen.csv};
\addlegendentry{coin framings}
\addplot[stwo, line width=1pt, densely dashed, mark=square*, mark size=1.5pt, mark options={solid}] table[col sep=comma, x=idx, y=log]{generated/stages-qwen.csv};
\addlegendentry{system-log framings}
\addlegendimage{muted, densely dotted}
\addlegendentry{fair coin}
\nextgroupplot[title={Llama 3.2\\3B}, width=1.5cm, xmin=-0.5, xmax=1.5, xtick={0,1}, xticklabels={base,Inst.}]
\addplot[muted, densely dotted, forget plot] coordinates {(-0.5,0.5) (1.5,0.5)};
\addplot[sone, line width=1pt, mark=*, mark size=1.7pt] table[col sep=comma, x=idx, y=fair]{generated/stages-llama.csv};
\addplot[stwo, line width=1pt, densely dashed, mark=square*, mark size=1.5pt, mark options={solid}] table[col sep=comma, x=idx, y=log]{generated/stages-llama.csv};
\nextgroupplot[title={OLMo 2\\1B}, width=3.2cm, xmin=-0.4, xmax=3.4, xtick={0,1,2,3}, xticklabels={base,SFT,DPO,RLVR}]
\addplot[muted, densely dotted, forget plot] coordinates {(-0.4,0.5) (3.4,0.5)};
\addplot[sone, line width=1pt, mark=*, mark size=1.7pt] table[col sep=comma, x=idx, y=fair]{generated/stages-olmo2.csv};
\addplot[stwo, line width=1pt, densely dashed, mark=square*, mark size=1.5pt, mark options={solid}] table[col sep=comma, x=idx, y=log]{generated/stages-olmo2.csv};
\nextgroupplot[title={OLMo 3 7B\\Instruct}, width=3.2cm, xmin=-0.4, xmax=3.4, xtick={0,1,2,3}, xticklabels={base,SFT,DPO,RLVR}]
\addplot[muted, densely dotted, forget plot] coordinates {(-0.4,0.5) (3.4,0.5)};
\addplot[sone, line width=1pt, mark=*, mark size=1.7pt] table[col sep=comma, x=idx, y=fair]{generated/stages-olmo3i.csv};
\addplot[stwo, line width=1pt, densely dashed, mark=square*, mark size=1.5pt, mark options={solid}] table[col sep=comma, x=idx, y=log]{generated/stages-olmo3i.csv};
\nextgroupplot[title={OLMo 3 7B\\Think}, width=3.2cm, xmin=-0.4, xmax=3.4, xtick={0,1,2,3}, xticklabels={base,SFT,DPO,RLVR}]
\addplot[muted, densely dotted, forget plot] coordinates {(-0.4,0.5) (3.4,0.5)};
\addplot[sone, line width=1pt, mark=*, mark size=1.7pt] table[col sep=comma, x=idx, y=fair]{generated/stages-olmo3t.csv};
\addplot[stwo, line width=1pt, densely dashed, mark=square*, mark size=1.5pt, mark options={solid}] table[col sep=comma, x=idx, y=log]{generated/stages-olmo3t.csv};
\end{groupplot}
\end{tikzpicture}

\vspace{2pt}
\ref{altlegend}
\caption{Probability that the next flip switches, averaged over flips 2 to 16 of the fair coins, when the letters are framed as coin flips (solid) or as a system log (dashed), along each training path. A model that treats the coin as fair sits on the dotted line. Where post-training raises switching, the log framing rises nearly as much; OLMo 3 never rises, and its Think path falls below 0.5 under both framings.}
\label{fig:switching}
\end{figure}

Post-training raised expected switching under coin framings for Qwen, from \stat{qwen_base/alt_fair} to \stat{qwen_instruct/alt_fair}, and along the OLMo 2 path, from \stat{olmo2_base/alt_fair} to \stat{olmo2_instruct/alt_fair} (Figure~\ref{fig:switching}). The system-log framing rose nearly as much: from \stat{qwen_base/alt_log} to \stat{qwen_instruct/alt_log} for Qwen and from \stat{olmo2_base/alt_log} to \stat{olmo2_instruct/alt_log} for OLMo 2. The coin-specific excess is small everywhere, at most \stat{llama_base/coin} (Llama base). H4 passed at its threshold: five post-trained checkpoints expect at least 0.52 switching under coin framings (Qwen Instruct, Llama Instruct and the three post-trained OLMo 2 checkpoints), and in four of them the log-framing excess is at least half the coin-framing excess. Llama Instruct misses narrowly, with a log excess of \stat{llama_instruct/log_excess} against a coin excess of \stat{llama_instruct/fair_excess}.

H6 failed. From OLMo 2 we predicted that OLMo 3's repetition avoidance would rise by at least 0.05 along both post-training paths and be mostly in place by DPO. Under the log framing it moved by \stat{h6/olmo3_instruct/rise} on the Instruct path and by \stat{h6/olmo3_think/rise} on the Think path. OLMo 3 hardly over-alternates at any stage. Under coin framings its Instruct checkpoints expect \stat{olmo3_instruct_sft/alt_fair} to \stat{olmo3_instruct/alt_fair}, and its Think checkpoints expect \stat{olmo3_think_sft/alt_fair} to \stat{olmo3_think/alt_fair}, below 0.5: they expect a fair coin to repeat more often than it does, under the log framing too (\stat{olmo3_think/alt_log}). The same Think checkpoints say that streaky coins are rigged (template \stat{olmo3_think/template}), the starkest gap between saying and expecting in the study.

\section{Reproducible is not real}
\label{sec:trap}

Before this study we ran a sealed reverse-correlation pilot in the style of superstitious-perception experiments \citep{gosselin2003superstitious}, whose first hypothesis predicted the confound described here. Qwen2.5 3B Instruct read 1,000 random $9 \times 9$ grids of \texttt{.} and \texttt{\#} characters and was asked, half the grids being said to hold ``a faint face'', whether each showed a hidden face or pure random static; the same question was asked for the letter X and for ``a zorp''. Scores were counterbalanced log-odds of the target against static, and classification images averaged the grids weighted by score. The images were highly reproducible across halves of the data for every question, the nonsense one included: split-half correlations were \stat{pilot/packed/face/split} (face) and \stat{pilot/packed/zorp/split} (zorp) for densely packed grids, and \stat{pilot/spaced/face/split} and \stat{pilot/spaced/zorp/split} for spaced grids, against a permutation-null 95th percentile of about \stat{pilot/packed/face/null}. The face image and the zorp image correlated at \stat{pilot/packed/cross} and \stat{pilot/spaced/cross}.

With exact log-probabilities the score is a deterministic function of the stimulus, so any question produces a reproducible map, and a permutation test against shuffled scores only confirms that the score depends on the stimulus. What these maps show is how the model reads the grid, not what it takes a face to be. We call this the deterministic-probe trap. A person asked to find a zorp has no systematic reason to answer differently for different grids, so a reproducible zorp map would be surprising. A model's answer is a fixed function of the grid whatever the question, so reproducibility is guaranteed and says nothing about the concept.

The coin study shows the same trap at a smaller scale. The nonsense twin of Qwen2.5 3B base correlates with switch rate at \stat{qwen_base/twin_template}, more strongly than its stated score does (\stat{qwen_base/template}), and the post-trained OLMo 2 1B checkpoints have stated templates between \stat{olmo2_instruct/template} and \stat{olmo2_dpo/template} at slopes of \stat{olmo2/max_slope} logits or less (Appendix Tables~\ref{tab:slopes} and~\ref{tab:readout}). A correlation between a score and a stimulus feature is not evidence of a judgment unless the score's size clearly exceeds that of a twin that cannot carry the concept. Our usability rule encodes that requirement, and the twin is what makes the agreement analysis in \S\ref{sec:agreement} interpretable.

\section{Discussion}

\paragraph{What post-training changes.} Post-trained models acquire a verbal judgment about randomness, that streaks look rigged, which matches what people write about randomness. Their likelihoods meanwhile keep or acquire separate habits: generic repetition avoidance in Qwen and OLMo 2, near-fair expectations in OLMo 3 Instruct, and an expectation of streaks in OLMo 3 Think. The verbal judgment is not read off the model's own predictions. This fits the finding that metalinguistic answers and direct probabilities can diverge \citep{hu2023prompting}, and it extends the generator--validator gap \citep{li2024benchmarking} to a setting where the model's ``generation'' is its predictive distribution over noise. Our data locate the change in post-training but do not show its mechanism.

\paragraph{Over-alternation needs a non-coin control.} \citet{vankoevering2024heads} found that over-alternation in generated flips is not explained by memorization or tokenization alone. In the likelihoods of the small open models studied here, it is almost entirely matched by the same letters framed as a log file. The measurements differ, generation against scoring and closed against open models, so the results do not conflict, but a claim that a model holds a coin-specific gambler's fallacy should be tested against such a framing.

\paragraph{A prediction that failed.} OLMo 2's rise in repetition avoidance did not carry over to OLMo 3, whose recipe and data differ in many ways. One family's trajectory through post-training is a weak guide to another's.

\paragraph{Base models.} Only OLMo 3's base checkpoint gave a usable answer in A/B format. That is a statement about the format, not evidence that base models lack a sense of randomness \citep{hu2024auxiliary}: their expectations can still be measured, and their stated scores, though too weak to use as judgments, tracked those expectations more closely than the post-trained models' answers did.

\section{Limitations}
\begin{itemize}
  \item The models are small (1B to 7B), from four families, on one task in one language with sequences of one length. Larger models may behave differently.
  \item Intervals come from resampling coins. They do not cover variation over wordings, of which there are only three, over prompt formats or over training seeds.
  \item Several verdicts sit at their bars (H4, H5 and one H8 pair). A different held-out seed could flip them.
  \item H8 was written after three of its five pairs were seen on calibration data, and the twin-partialled check and the upturn diagnostic were added after the evaluation run.
  \item The registered curve criterion is lenient (\S\ref{sec:curves}). The comparison with human judgments is qualitative: we collected no human ratings, and two exploratory items, a case-alternated spelling control and a matched comparison with human ratings, were not run.
  \item Base and post-trained checkpoints also differ in prompt format on stated questions (raw text \texttt{\textbackslash nAnswer:} against chat templates), whereas implicit scores always use raw text.
  \item Stated wordings \texttt{w1} (``rigged or patterned'') and \texttt{w3} (``non-random'') cue streaks and yield 3--6$\times$ steeper slopes than \texttt{w2} (``a trick coin''; Appendix Table~\ref{tab:slopes}), while the implicit framings (``rigged'', ``trick'', ``loaded'') pick up heads--tails bias more than streaks (Appendix Table~\ref{tab:sensitivity}).
  \item Neither score was counterbalanced by swapping \texttt{H} and \texttt{T}. Partialling out flip~1 and heads share removes the linear part post hoc, not nonlinear bias.
  \item The Think checkpoints put less probability on A or B (median \stat{olmo3_think_dpo/mass} to \stat{olmo3_think/mass}), so their stated scores depend on closing the thought empty rather than letting the model reason.
  \item All scoring ran on one hardware and software stack. We checked the shared-prefix computation for exactness there but did not replicate on other hardware.
\end{itemize}

\section{Reproducibility}
The protocol (SHA-256 \texttt{\stat{sha12}\ldots}, sealed \stat{sealed}), the amendments file, the analysis code with its pinned lock file, the raw per-coin scores for every checkpoint and the earlier pilot record are released with the paper. Rerunning the evaluation takes about six hours on one laptop. Every result, table and figure in this paper is generated from the raw scores and the sealed protocol by one command.

\bibliographystyle{plainnat}
\bibliography{refs}

\appendix

\section{Prompts}
\label{app:prompts}

A stated prompt under wording w1, with the rigged option listed first, reads:

\begin{quote}\small\ttfamily
Here are 16 coin flips (H = heads, T = tails):\\
H T T H H H T H T T H T H H T T\\[4pt]
Half of sequences come from a fair random coin; half are rigged or patterned.\\
A) a rigged or patterned sequence\\
B) fair random coin flips\\
Reply with a single letter (A or B).
\end{quote}

\noindent The other wordings change the lead, the question and the options:

\begin{description}[font=\normalfont\itshape, leftmargin=1.2em]
  \item[w2] ``A friend flipped a coin 16 times and wrote down the results (H = heads, T = tails):'' / ``Some friends use a fair coin and some use a trick coin, equally often.'' / ``a trick coin'' or ``a fair coin''.
  \item[w3] ``Sequence of 16 coin tosses, H for heads and T for tails:'' / ``Was this produced by a genuinely random process or by a non-random one? Both are equally common.'' / ``a non-random process'' or ``a genuinely random process''.
\end{description}

\noindent The nonsense twins keep each lead and replace the question and options with ``Half of sequences are zorps; half are blargs.'' (``a zorp sequence'' or ``a blarg sequence''), ``Some friends use a zorp coin and some use a blarg coin, equally often.'' (``a zorp coin'' or ``a blarg coin'') and ``Was this produced by a zorp process or by a blarg one? Both are equally common.'' (``a zorp process'' or ``a blarg process'').

\noindent The implicit framings, each followed by a newline and the flips, are ``Here are 16 flips of a fair coin (H = heads, T = tails):'' against ``\ldots of a rigged coin \ldots''; ``I flipped a fair coin 16 times and got:'' against ``\ldots a trick coin \ldots''; and ``Results of 16 tosses of an unbiased coin (H/T):'' against ``\ldots of a loaded coin \ldots''. The log framings are ``Here is a 16-character string from a system log (each character is H or T):'' and ``Status codes from 16 consecutive readings of a machine (each is H or T):''.

\section{Models}
\label{app:models}

\begin{table}[H]
\centering
\footnotesize
\caption{Checkpoints, Hugging Face repositories and pinned revisions. Prompt mode applies to stated questions; implicit scores always use raw text.}
\label{tab:models}
\begin{tabular}{@{}llll@{}}
\toprule
Checkpoint & Repository & Revision & Prompt mode \\
\midrule
\rows{generated/models}
\bottomrule
\end{tabular}
\end{table}

\section{Additional tables}

\begin{table}[H]
\centering
\footnotesize
\caption{Slopes of the stated score and of its nonsense twin on switch rate, in logits per unit switch rate, for each wording. An asterisk marks a wording that passes the usability rule.}
\label{tab:slopes}
\begin{tabular}{@{}lccccccc@{}}
\toprule
 & \multicolumn{2}{c}{w1} & \multicolumn{2}{c}{w2} & \multicolumn{2}{c}{w3} & \\
\cmidrule(lr){2-3}\cmidrule(lr){4-5}\cmidrule(lr){6-7}
Checkpoint & stated & twin & stated & twin & stated & twin & Usable \\
\midrule
\rows{generated/slopes}
\bottomrule
\end{tabular}
\end{table}

\begin{table}[H]
\centering
\footnotesize
\caption{Readout controls across fair coins (post hoc). Templates of the stated score and of the nonsense twin; correlations among stated, twin and implicit scores; and the stated--implicit correlation with the twin partialled out.}
\label{tab:readout}
\begin{tabular}{@{}lcccccc@{}}
\toprule
Checkpoint & \makecell{Stated\\template} & \makecell{Twin\\template} & \makecell{$r$(stated,\\twin)} & \makecell{$r$(stated,\\implicit)} & \makecell{$r$(twin,\\implicit)} & \makecell{Partial\\$r$} \\
\midrule
\rows{generated/readout}
\bottomrule
\end{tabular}
\end{table}

\begin{table}[H]
\centering
\footnotesize
\caption{Sensitivity (exploratory): area under the ROC curve for each rigged set against fair coins, using the stated score (averaged over wordings) and the implicit score (averaged over framings). Values above 0.5 mean the rigged set scores as more rigged.}
\label{tab:sensitivity}
\begin{tabular}{@{}lcccccc@{}}
\toprule
 & \multicolumn{2}{c}{Streaky} & \multicolumn{2}{c}{Alternating} & \multicolumn{2}{c}{Biased} \\
\cmidrule(lr){2-3}\cmidrule(lr){4-5}\cmidrule(lr){6-7}
Checkpoint & stated & implicit & stated & implicit & stated & implicit \\
\midrule
\rows{generated/sensitivity}
\bottomrule
\end{tabular}
\end{table}

\begin{table}[H]
\centering
\footnotesize
\caption{Mean pooled stated score by number of switches $k$ (Figure~\ref{fig:curves}).}
\label{tab:curves}
\pgfplotstabletypeset[
  col sep=comma,
  columns={rate,qweninstruct,olmo3instruct,olmo3think},
  columns/rate/.style={column name={Switch rate}, fixed, precision=2, zerofill},
  columns/qweninstruct/.style={column name={Qwen2.5 3B Instruct}, fixed, precision=2, zerofill},
  columns/olmo3instruct/.style={column name={OLMo 3 7B Instruct}, fixed, precision=2, zerofill},
  columns/olmo3think/.style={column name={OLMo 3 7B Think}, fixed, precision=2, zerofill},
  every head row/.style={before row=\toprule, after row=\midrule},
  every last row/.style={after row=\bottomrule},
]{generated/curves.csv}
\end{table}

\begin{table}[H]
\centering
\footnotesize
\caption{Feature correlations and H/T controls across fair coins (post hoc). Correlations with flip~1 being \texttt{H} (stated, twin, implicit), heads share $h$ (stated), imbalance $|2h - 1|$ (implicit), and stated--implicit agreement after partialling out flip~1 ($r \mid \mathrm{first\ H}$) or flip~1 and heads share ($r \mid \mathrm{H/T}$).}
\label{tab:features}
\setlength{\tabcolsep}{4pt}
\begin{tabular}{@{}lccccccc@{}}
\toprule
Checkpoint & \makecell{Says\\first H} & \makecell{Twin\\first H} & \makecell{Expects\\first H} & \makecell{Says\\heads} & \makecell{Expects\\imb.} & \makecell{$r \mid$\\first H} & \makecell{$r \mid$\\H/T} \\
\midrule
\rows{generated/features}
\bottomrule
\end{tabular}
\end{table}

\begin{table}[H]
\centering
\footnotesize
\caption{Base-to-post-trained agreement shifts after partialling out flip~1 and heads share (post hoc), with paired bootstrap 95\% intervals. Levels are in Table~\ref{tab:features}.}
\label{tab:symmetry}
\setlength{\tabcolsep}{4pt}
\begin{tabular}{@{}lcccc@{}}
\toprule
Pair & \makecell{Shift $\mid$ first H\\{[95\% CI]}} & \makecell{H/T\\base $\to$ post} & \makecell{Shift $\mid$ H/T\\{[95\% CI]}} & \makecell{Shift $\mid$ H/T, twin\\{[95\% CI]}} \\
\midrule
\rows{generated/symmetry}
\bottomrule
\end{tabular}
\end{table}

\section{Amendments after the evaluation}
\label{app:amendments}
\begin{enumerate}
  \item \emph{Curve criterion.} The registered criterion also accepts curves that fall from the streaky end and then flatten. H7 stays a pass as registered; the upturn at perfect alternation is reported as a post-hoc diagnostic.
  \item \emph{Disclosure.} Three of the five H8 pairs had been seen on calibration data before H8 was written. Only the two OLMo 3 pairs were predictions.
  \item \emph{Not run.} The case-alternated spelling control and the matched comparison with human ratings, both listed as exploratory, were not run on the evaluation seed.
  \item \emph{First-flip bias.} Flip~1 and heads share explain most of the stated--twin correlation and drive the apparent H8 drop in Llama and OLMo 2. Partialling them out keeps the drop in Qwen and both OLMo 3 paths; H8 stays a pass as registered.
  \item \emph{Follow-up scripts.} Added unsealed scripts for wording-by-framing tables, raw-text runs on post-trained models, open-thought runs on OLMo 3 Think, and larger base/instruct pairs. No sealed file was changed.
\end{enumerate}

\end{document}
